% =====================================================================
%  Выпускная квалификационная работа
%  «Адаптивные топологии мульти-агентных LLM-систем
%   с включением человека в контур управления»
%
%  Образовательная программа «Прикладной анализ данных и
%  искусственный интеллект», НИУ ВШЭ — Санкт-Петербург, 2026
%
%  Сборка:
%    pdflatex main.tex
%    pdflatex main.tex     (второй прогон для оглавления/ссылок)
%
%  Если установлен latexmk:
%    latexmk -pdf main.tex
% =====================================================================

\documentclass[14pt,a4paper]{extarticle}

%%% --- Кодировка и язык --------------------------------------------- %%%
\usepackage[T2A]{fontenc}
\usepackage[utf8]{inputenc}
\usepackage[english,russian]{babel}

%%% --- Шрифт с полной поддержкой кириллицы ------------------------- %%%
% Linux Libertine --- бесплатный шрифт с засечками, visually
% близок к Times New Roman, имеет тонкое начертание (меньше
% x-height, чем у PT Serif), полная T2A кириллица.
% Методичка ВКР НИУ ВШЭ допускает отступления от TNR
% (Приложение 6 п. 2.2 --- «в идеале TNR, отступления возможны»).
\usepackage{libertine}
\usepackage[T2A,T1]{fontenc}
\usepackage{microtype}

%%% --- Поля, междустрочный интервал, отступ -------------------------- %%%
% Методичка: левое поле >= 2см. Стандарт ВКР: левое 2.5 cm, правое 1.5 cm.
\usepackage[left=2.5cm, right=1.5cm, top=2cm, bottom=2cm]{geometry}
\usepackage{setspace}
\setstretch{1.3}                 % 1.3 (в пределах допуска методички 1.3--1.5)
\usepackage{indentfirst}
\setlength{\parindent}{1cm}      % красная строка 1 см
\setcounter{tocdepth}{1}         % в оглавлении --- только до subsection

%%% --- Графика, ссылки, таблицы, формулы ---------------------------- %%%
\usepackage{graphicx}
\usepackage{xcolor}
\usepackage{tikz}
\usetikzlibrary{positioning, shapes.geometric, arrows.meta, fit,
                backgrounds, calc}
\usepackage[hidelinks]{hyperref}
\usepackage{csquotes}
\usepackage{enumitem}
\setlist{nosep, leftmargin=*, labelindent=\parindent}
\usepackage{amsmath,amssymb,amsthm}
\usepackage{booktabs}
\usepackage{array}
\usepackage{tabularx}
\usepackage{longtable}
\usepackage{caption}
\captionsetup[figure]{labelfont=bf,format=plain,justification=centering,
                      labelsep=period}
\captionsetup[table]{labelfont=bf,format=plain,justification=raggedright,
                     singlelinecheck=false,labelsep=period}

%%% --- Подписи на русском ----------------------------------------- %%%
\addto\captionsrussian{%
  \renewcommand{\figurename}{Рисунок}%
  \renewcommand{\tablename}{Таблица}%
  \renewcommand{\contentsname}{Оглавление}%
  \renewcommand{\refname}{Библиографический список}%
  \renewcommand{\appendixname}{Приложение}%
}

%%% --- Нумерация рисунков/таблиц с указанием раздела ---------------- %%%
\numberwithin{figure}{section}
\numberwithin{table}{section}
\numberwithin{equation}{section}

%%% --- Стиль заголовков -------------------------------------------- %%%
\usepackage{titlesec}
\titleformat{\section}
    {\normalfont\Large\bfseries}{\thesection}{1em}{}
\titleformat{\subsection}
    {\normalfont\large\bfseries}{\thesubsection}{1em}{}
\titleformat{\subsubsection}
    {\normalfont\normalsize\bfseries}{\thesubsubsection}{1em}{}
\titlespacing*{\section}{0pt}{0.8\baselineskip}{0.4\baselineskip}
\titlespacing*{\subsection}{0pt}{0.6\baselineskip}{0.3\baselineskip}
\titlespacing*{\subsubsection}{0pt}{0.5\baselineskip}{0.25\baselineskip}

%%% --- Колонтитулы ------------------------------------------------- %%%
\pagestyle{plain}

%%% --- Листинги (для YAML-конфигов в приложении Б) ----------------- %%%
\usepackage{listings}
\usepackage{multicol}
\usepackage[section]{placeins}      % \FloatBarrier на границах разделов
\renewcommand{\textfraction}{0.05}  % минимум 5% текста на странице с float
\renewcommand{\topfraction}{0.95}   % float может занимать до 95% верха
\renewcommand{\bottomfraction}{0.95}
\renewcommand{\floatpagefraction}{0.6} % float-only page если >=60% заполнено
\lstset{
    basicstyle=\ttfamily\small,
    breaklines=true,
    columns=fullflexible,
    keepspaces=true,
    showstringspaces=false,
    frame=single,
    rulecolor=\color{gray!60},
    framerule=0.4pt,
}

%%% --- Гиперссылки в оглавлении ----------------------------------- %%%
\hypersetup{
    colorlinks=true,
    linkcolor=black,
    citecolor=black,
    urlcolor=blue,
    linktoc=all,
    bookmarksnumbered=true
}

%%% --- Глобальные команды ----------------------------------------- %%%
\newcommand{\code}[1]{\texttt{#1}}
\newcommand{\TODO}[1]{{\color{red}\textbf{[TODO]}~#1}}
\newcommand{\PLACEHOLDER}[1]{\textcolor{orange}{[#1]}}

% =====================================================================
\begin{document}

%%% ===== Титульный лист ===== %%%
% Титульный лист по образцу Приложения 5 методички ВКР.
% Имя научного руководителя пользователь вписывает сам.

\thispagestyle{empty}
\begin{center}
\large

ФЕДЕРАЛЬНОЕ ГОСУДАРСТВЕННОЕ АВТОНОМНОЕ ОБРАЗОВАТЕЛЬНОЕ\\
УЧРЕЖДЕНИЕ ВЫСШЕГО ОБРАЗОВАНИЯ\\
\textbf{«НАЦИОНАЛЬНЫЙ ИССЛЕДОВАТЕЛЬСКИЙ УНИВЕРСИТЕТ\\
«ВЫСШАЯ ШКОЛА ЭКОНОМИКИ»}

\vspace{0.7cm}
Факультет «Школа информатики, физики и технологий»

\vfill

Соснин Тимофей Алексеевич

\vspace{0.8cm}
\Large
\textbf{Адаптивные топологии мульти-агентных\\
систем больших языковых моделей с включением\\
человека в контур управления}

\vspace{0.4cm}
\large
Выпускная квалификационная работа\\
\textsc{бакалаврская работа}

\vspace{0.6cm}
по направлению подготовки 01.03.02\\
«Прикладная математика и информатика»\\[0.2cm]
образовательная программа\\
«Прикладной анализ данных и искусственный интеллект»

\vfill

\begin{flushright}
\begin{tabular}{p{6.5cm}}
Руководитель \\[0.1cm]
к.т.н., доцент департамента информатики \\[0.4cm]
\rule{4.5cm}{0.3pt} \\
Д.А. Рюмин \\
\end{tabular}
\end{flushright}

\vfill
Санкт-Петербург, 2026
\end{center}


%%% ===== Оглавление ===== %%%
\newpage
\tableofcontents
\thispagestyle{empty}

%%% ===== Аннотация RU + EN ===== %%%
\newpage
% Аннотация на русском и английском, ~150 слов каждая.
% Заполняется в финале после получения результатов экспериментов.

\section*{Аннотация}
\addcontentsline{toc}{section}{Аннотация}

\TODO{Финал. Заполняется после Главы 5. ~150 слов на русском.
Краткое описание предметной области, цели, методологии,
ключевых результатов и вклада работы.}

\vspace{1cm}

\selectlanguage{english}
\section*{Abstract}
\addcontentsline{toc}{section}{Abstract}

\TODO{Final. Filled after Chapter 5. About 150 words in English.}

\selectlanguage{russian}


%%% ===== Ключевые слова (без \newpage --- сразу после аннотации) ===== %%%
% Список ключевых слов, 5--10 единиц.

\section*{Ключевые слова}
\addcontentsline{toc}{section}{Ключевые слова}

\noindent мульти-агентные системы, большие языковые модели,
коммуникационная топология, адаптивная маршрутизация, человек
в~контуре управления, оракул адаптации, кросс-семейная валидация,
Парето-оптимизация, бутстрап-оценивание, оркестрация
LLM-агентов.

\vspace{0.5cm}

\selectlanguage{english}
\noindent\textbf{Keywords:} multi-agent systems, large language
models, communication topology, adaptive routing, human-in-the-loop,
adaptation oracle, cross-family validation, Pareto optimization,
bootstrap inference, LLM agent orchestration.
\selectlanguage{russian}


%%% ===== Введение ===== %%%
\newpage
% =====================================================================
% Введение.
% Строгий порядок по правкам друга-ревьюера
%   1. Актуальность работы
%   2. Новизна
%   3. Предмет исследования
%   4. Объект исследования
%   5. Цель работы (без слова «оценка», только «количественная оценка»)
%   6. Аналитический обзор предлагаемых решений (краткий)
%   7. Практическая значимость
%   8. Теоретическая значимость
%   9. Методы исследования
%  10. Задачи работы
%  11. Структура работы (с числом таблиц / рисунков / страниц / формул)
%
% Объем 3--4 страницы.
% =====================================================================

\section*{Введение}
\addcontentsline{toc}{section}{Введение}

\textbf{Актуальность работы.} В 2023--2026 годах применение больших
языковых моделей (от англ. Large Language Models, LLM) перешло от
одиночных диалоговых сценариев к кооперативной работе нескольких
автономных агентов, образующих мульти-агентную систему
(от англ. Multi-Agent System, MAS). Открытые фреймворки (AutoGen,
ChatDev, MetaGPT, Magentic-One) показывают, что разделение труда
между специализированными агентами повышает качество решения
сложных задач, но каждый из них жестко фиксирует структуру
коммуникации на этапе проектирования; систематических данных об
оптимальной топологии для каждого класса задач в открытой литературе
нет. Параллельный интерес к интеграции человека в контур управления
ограничивается единичными паттернами без сравнения эффективности.
Сочетание этих обстоятельств определяет научную и прикладную
актуальность работы.

\textbf{Научная новизна.} В настоящей работе впервые проведено
прямое сравнение пяти базовых коммуникационных топологий
(звезда, цепочка, полносвязная, дебатная, иерархическая) на
едином программном стенде, на одной языковой модели и на одной
выборке задач из четырех различных классов. Предложена и
реализована мета-топология, переключающая активную базовую
конфигурацию в зависимости от текущей фазы решения и наблюдаемых
сигналов агентов. Введена таксономия из пяти специфичных для
адаптивных мульти-агентных систем метрик ---
$N_{\text{switch}}$, $r_{\text{guard}}$, $\gamma$,
$r_{\text{hurt}}$, $\Delta_{\text{oracle}}$, --- позволяющая
количественно характеризовать поведение мета-топологии. Впервые
сформулирована и экспериментально проверена матрица совместимости
пяти ролей человека в контуре управления с пятью базовыми
топологиями.

\textbf{Предмет исследования.} Коммуникационная топология
мульти-агентной системы больших языковых моделей, механизм
ее динамической адаптации в процессе решения задачи, а также
позиция человека в графе обмена сообщениями между агентами.

\textbf{Объект исследования.} Мульти-агентные системы на основе
больших языковых моделей, решающие сложные задачи различных
классов в кооперации со специализированными агентами и
участником-человеком.

\textbf{Цель работы.} Количественная оценка влияния выбора
коммуникационной топологии мульти-агентной системы, механизма
ее адаптации и позиции человека в контуре управления на
эффективность решения задач разных классов с одновременным
учетом качества решения, его стоимости и времени выполнения.

\textbf{Аналитический обзор предлагаемых решений.} В 2022--2026
годах выделяются три направления развития MAS LLM. Архитектурные
фреймворки (AutoGen, ChatDev, MetaGPT) фиксируют топологию и не
сравнивают альтернативы. Автоматическое проектирование графов
(G-Designer, GPTSwarm) подбирает структуру однократно и не
пересматривает ее. Динамическая адаптация (DyLAN, AgentVerse)
меняет состав команды, но не структуру связей. Вопрос роли человека
освещен фрагментарно, без количественной оценки ролей. Подробный
анализ --- в главе~\ref{sec:review}, сравнительная
таблица~\ref{tab:rev:methods} восемнадцати методов по шести
критериям.

\textbf{Практическая значимость.} Работа дает разработчикам MAS
количественные данные о том, какая топология эффективнее для
каждого из четырех рассмотренных классов задач и когда имеет
смысл адаптивное переключение. Программная инфраструктура
реализована как открытый фреймворк адаптивных топологий
(от англ. Adaptive Topologies for Multi-agent systems, ATM),
исходный код --- в репозитории GitHub (приложение~\ref{app:github}),
допускает повторное использование и расширение. Сформулированные
правила выбора топологии и роли человека под класс задачи применимы
при проектировании интеллектуальных ассистентов в программировании,
аналитике и принятии решений.

\textbf{Теоретическая значимость.} Введенная таксономия
адаптивно-специфичных метрик ($N_{\text{switch}}$, $r_{\text{guard}}$,
$\gamma$, $r_{\text{hurt}}$, $\Delta_{\text{oracle}}$) позволяет
количественно характеризовать поведение динамически перестраиваемых
MAS по единой шкале. Формализация MAS как кортежа из шести
компонентов с фазово-зависимой коммуникационной структурой
расширяет теоретическую базу; эмпирическая проверка гипотезы о
фазовой специфике топологий дает обоснование для дальнейших
теоретических работ.

\textbf{Методы исследования.} Комплекс взаимодополняющих методов:
аналитический обзор литературы по ключевым словам в открытых базах;
формализация MAS средствами теории графов и конечных автоматов;
программная реализация на Python с LangGraph; контролируемый
эксперимент с предварительной регистрацией гипотез и порогов
улучшения; статистическая обработка (дисперсионный анализ от англ.
Analysis of Variance, ANOVA; парные сравнения с поправкой
Бенджамини--Хохберга на уровне ложных открытий от англ. False
Discovery Rate, FDR; коэффициент Коэна; непараметрический
бутстрап с доверительными интервалами от англ. Confidence
Interval, CI); оракульный референс per-task best-of
для верхней границы качества адаптации; имитационное моделирование
участника-человека языковой моделью с подсказкой по роли.

\textbf{Задачи работы.} Для достижения сформулированной выше
цели в работе решаются следующие задачи.
\begin{enumerate}
  \item Провести аналитический обзор работ 2022--2026 годов по
        мульти-агентным системам больших языковых моделей,
        коммуникационным топологиям и интеграции человека в
        контур управления, выделить исследовательскую лакуну.
  \item Сформулировать формальную модель мульти-агентной системы
        как кортежа из множества агентов, множества ребер
        коммуникации, множества ролей агентов и человека,
        задачи и множества фаз решения.
  \item Разработать и реализовать программный фреймворк, в
        котором пять базовых топологий и адаптивная мета-топология
        реализованы в единой архитектуре на основе библиотеки
        LangGraph.
  \item Спроектировать систему оценочных метрик, включающую
        базовые показатели качества, стоимости и времени, а также
        специфические для адаптивной мета-топологии характеристики
        переключений, заблокированных решений маршрутизатора,
        стоимостного вклада маршрутизатора, доли регрессий
        качества и разрыва до верхнего потолка адаптации.
  \item Провести воронку из четырех экспериментов и
        подтверждающего прогона на четырех классах задач ---
        программирование, математические рассуждения, творческая
        генерация, анализ данных --- с предварительно
        зафиксированными гипотезами и порогами улучшения.
  \item Выполнить статистический анализ результатов и сформулировать
        ответы на четыре исследовательских вопроса (от англ.
        Research Question, RQ), описанных в
        разделе~\ref{sec:meth:funnel}.
  \item Выработать практические рекомендации по выбору топологии
        и роли человека в зависимости от класса задачи.
\end{enumerate}

\textbf{Структура работы.} Работа состоит из введения, пяти глав
основного содержания, раздела авторского вклада, заключения,
библиографического списка и четырех приложений (А, Б, В, Г).
Глава~1 --- аналитический обзор литературы и исследовательская
лакуна. Глава~2 --- формальная модель, шесть топологий, пять
ролей человека, метрики, корпус, дизайн воронки. Глава~3 ---
архитектура программной системы. Глава~4 --- результаты четырех
основных экспериментов и подтверждающих прогонов на дополнительном
семействе моделей. Глава~5 --- анализ, ответы на четыре RQ и
угрозы валидности. Раздел авторского вклада фиксирует оригинальные
элементы; заключение подводит итоги и формулирует направления
дальнейших исследований. Приложения --- функциональная схема,
глоссарий, состав репозитория, ссылка на исходный код.

В работе пять глав основного содержания, четыре приложения,
19~таблиц, 8~рисунков и~5~формул; общий объем работы с~введением,
заключением, библиографическим списком и~приложениями составляет
62~страницы; библиографический список --- 44~источника.


%%% ===== Главы ===== %%%
\newpage
% =====================================================================
% Глава 1. Аналитический обзор литературы.
% Объем 6--7 страниц.
% Структура по правкам друга-ревьюера
%   1.1 Мульти-агентные системы больших языковых моделей
%   1.2 Коммуникационные топологии агентов
%   1.3 Механизмы адаптации топологий
%   1.4 Безопасность и устойчивость топологий
%   1.5 Человек в контуре управления мульти-агентными системами
%   1.6 Сравнительная аналитическая таблица методов
%   1.7 Что отсутствует в существующих работах
%   1.8 Выводы по главе
% =====================================================================

\section{Аналитический обзор литературы}
\label{sec:review}

Настоящая глава представляет аналитический обзор работ, на стыке
которых формируется предметная область исследования. Обзор
охватывает пять направлений --- архитектуру мульти-агентных систем
на основе больших языковых моделей, коммуникационные топологии
агентов, механизмы их адаптации, безопасность и устойчивость
топологий, а также интеграцию человека в контур управления
мульти-агентными системами. Глава завершается сравнительной
таблицей основных методов и определением исследовательской лакуны,
заполнение которой составляет содержательный вклад настоящей
работы.

Методология обзора. В рассматриваемую выборку включены работы
с 2022 по 2026 годы, отобранные по ключевым словам \emph{multi-agent
LLM}, \emph{agent topology}, \emph{LLM debate}, \emph{human-in-the-
loop multi-agent} в базах публикаций arXiv, ACL Anthology, ICLR,
NeurIPS, ICML. Из первоначального списка из более чем шестидесяти
работ оставлены те, которые либо предлагают новую архитектуру
взаимодействия агентов, либо систематически анализируют свойства
существующих архитектур. Работы, посвященные исключительно
обучению одиночных моделей или классическому планированию в
дискретных средах, в обзор не включены.

\subsection{Мульти-агентные системы больших языковых моделей}
\label{sec:rev:mas}

Идея использования нескольких автономных программных агентов для
совместного решения задач восходит к классическим работам по
распределенному искусственному интеллекту. С появлением больших
языковых моделей (от англ. Large Language Models, LLM)
концепция мульти-агентных систем (от англ. Multi-Agent System,
MAS) переживает второе рождение --- роль агента теперь играет
языковая модель с заданной системной подсказкой и доступным
набором инструментов. Обзорная работа~\cite{wang2024} систематизирует
быстрорастущую совокупность исследований и выделяет несколько
устойчивых архитектурных образцов.

Первой широко известной системой подобного класса является
фреймворк AutoGen~\cite{wu2023}, в котором несколько агентов
обмениваются сообщениями в гибком формате многораундового диалога
и могут вызывать внешние инструменты, в том числе подключать
человека. AutoGen не фиксирует топологию заранее --- структура
взаимодействия определяется императивным кодом конкретного
сценария. Это позволяет реализовать любую конфигурацию, но
лишает фреймворк возможности сравнивать топологии между собой
в рамках единой архитектуры.

Иной образец продемонстрирован в системе ChatDev~\cite{qian2024},
где процесс разработки программного обеспечения моделируется
последовательностью ролевых агентов --- генерального директора,
программиста, тестировщика --- работающих по принципу
водопадной модели. Метод MetaGPT~\cite{hong2024} развивает
идею за счет введения стандартных операционных процедур, что
позволяет снизить число ошибок согласования между ролями.
Эти работы демонстрируют, что разделение труда между специализированными
агентами повышает качество решения сложных задач, однако
топология взаимодействия в них жестко зашита в архитектуру.

Параллельно развивается направление, в котором отдельный агент
итеративно улучшает собственный ответ. Метод Self-Refine
\cite{selfrefine2024} реализует цикл генерация --- критика ---
переработка, а метод Reflexion~\cite{reflexion2023} расширяет
цикл вербальным подкреплением --- модель формирует словесную
обратную связь и хранит ее в эпизодической памяти. Несмотря
на привлекательную простоту, эти методы по существу остаются
одноагентными и не используют разнообразие точек зрения,
доступное мульти-агентным конфигурациям.

\subsection{Коммуникационные топологии агентов}
\label{sec:rev:topology}

Под коммуникационной топологией понимается граф, описывающий,
какие агенты могут обмениваться сообщениями и в каком порядке.
В существующих фреймворках топология выбирается архитектурно
и не пересматривается во время выполнения задачи. Тем не менее
ряд недавних работ ставит вопрос об оптимальной структуре
такого графа. Используемая в настоящей работе таксономия графовых
структур включает пять базовых типов --- звезда (центральный
координатор и периферийные работники), цепочка (линейный конвейер),
полносвязная (каждый агент общается с каждым), дебатная (два
оппонента и арбитр) и иерархическая (двух- или многоуровневое
разделение труда). Указанные пять типов и составляют пространство
сравнения, исследуемое в настоящей работе. Схематическое
изображение пяти базовых топологий приведено на
рисунке~\ref{fig:rev:topologies}.

\begin{figure}[!htbp]
\centering
\begin{tikzpicture}[
  font=\scriptsize,
  node distance=0.35cm,
  every node/.style={align=center},
  ag/.style={circle, draw=black, thick, minimum size=4.5mm,
             inner sep=0.5pt, fill=blue!10},
  hub/.style={circle, draw=black, thick, minimum size=5mm,
              inner sep=0.5pt, fill=orange!25},
  jud/.style={circle, draw=black, thick, minimum size=5mm,
              inner sep=0.5pt, fill=green!18},
  ed/.style={-{Latex[length=1mm]}, thick},
  bid/.style={{Latex[length=1mm]}-{Latex[length=1mm]}, thick},
]

% --- Star ---
\begin{scope}[xshift=0cm]
  \node[hub] (s_c) at (0,0) {C};
  \node[ag] (s_1) at (90:0.95)  {1};
  \node[ag] (s_2) at (162:0.95) {2};
  \node[ag] (s_3) at (234:0.95) {3};
  \node[ag] (s_4) at (306:0.95) {4};
  \node[ag] (s_5) at (18:0.95)  {5};
  \foreach \i in {s_1,s_2,s_3,s_4,s_5}{ \draw[bid] (s_c) -- (\i); }
  \node[below=1.2cm of s_c, font=\footnotesize\bfseries] {Star};
\end{scope}

% --- Chain ---
\begin{scope}[xshift=3.0cm]
  \node[ag] (c_1) at (-1.0,0) {1};
  \node[ag] (c_2) at (-0.0,0) {2};
  \node[ag] (c_3) at (1.0,0)  {3};
  \draw[ed] (c_1) -- (c_2);
  \draw[ed] (c_2) -- (c_3);
  \node[below=1.1cm of c_2, font=\footnotesize\bfseries] {Chain};
\end{scope}

% --- Mesh ---
\begin{scope}[xshift=6.5cm]
  \node[ag] (m_1) at (90:0.85)  {1};
  \node[ag] (m_2) at (18:0.85)  {2};
  \node[ag] (m_3) at (306:0.85) {3};
  \node[ag] (m_4) at (234:0.85) {4};
  \node[ag] (m_5) at (162:0.85) {5};
  \foreach \i/\j in {m_1/m_2, m_1/m_3, m_1/m_4, m_1/m_5,
                     m_2/m_3, m_2/m_4, m_2/m_5,
                     m_3/m_4, m_3/m_5, m_4/m_5} {
    \draw[thick, gray!70] (\i) -- (\j);
  }
  \node[below=1.1cm of m_3, font=\footnotesize\bfseries] {Mesh};
\end{scope}

% --- Debate (вторая строка) ---
\begin{scope}[xshift=1.5cm, yshift=-3.4cm]
  \node[ag] (d_p) at (-0.8, 0.4) {$+$};
  \node[ag] (d_n) at (0.8, 0.4)  {$-$};
  \node[jud] (d_j) at (0, -0.7)  {J};
  \draw[bid] (d_p) -- (d_n);
  \draw[ed] (d_p) -- (d_j);
  \draw[ed] (d_n) -- (d_j);
  \node[below=0.4cm of d_j, font=\footnotesize\bfseries] {Debate};
\end{scope}

% --- Hierarchical ---
\begin{scope}[xshift=5.5cm, yshift=-3.4cm]
  \node[hub] (h_c) at (0, 0.6) {C};
  \node[ag] (h_a1) at (-1.0, -0.2) {a1};
  \node[ag] (h_a2) at (-0.3, -0.2) {a2};
  \node[ag] (h_b1) at (0.3, -0.2)  {b1};
  \node[ag] (h_b2) at (1.0, -0.2)  {b2};
  \draw[ed] (h_c) -- (h_a1);
  \draw[ed] (h_c) -- (h_a2);
  \draw[ed] (h_c) -- (h_b1);
  \draw[ed] (h_c) -- (h_b2);
  \draw[bid, gray!70] (h_a1) -- (h_a2);
  \draw[bid, gray!70] (h_b1) -- (h_b2);
  \node[below=0.85cm of h_c, font=\footnotesize\bfseries]
       {Hierarchical};
\end{scope}

\end{tikzpicture}
\caption{Схематическое изображение пяти базовых коммуникационных
топологий, рассматриваемых в работе. Узлы --- агенты; C ---
координатор, J --- арбитр, $+$ и $-$ --- оппоненты в дебатах,
a1, a2, b1, b2 --- участники подкоманд. Сплошные направленные
ребра обозначают передачу управления, серые двунаправленные ---
обмен сообщениями в общей шине.}
\label{fig:rev:topologies}
\end{figure}

Работа GPTSwarm~\cite{zhuge2024} рассматривает мульти-агентную
систему как обучаемый граф вычислений и оптимизирует его
параметры на корпусе задач. Авторы показывают, что
оптимизированные графы превосходят случайные конфигурации,
но в качестве пространства поиска используют только статичные
структуры. Метод G-Designer~\cite{zhang2024} применяет
графовые нейронные сети (от англ. Graph Neural Network ---
GNN) для автоматического подбора топологии под конкретную
задачу. Несмотря на адаптацию к типу задачи, выбор топологии
производится однократно и не пересматривается в ходе
решения.

В работе~\cite{qian2024scaling} систематически исследуются
закономерности масштабирования мульти-агентных систем при
увеличении числа агентов. Авторы обнаруживают, что топологии
типа малого мира (small-world) демонстрируют повышенную
робастность к ошибкам отдельных агентов, однако сравнение
ведется между статичными конфигурациями и не учитывает
динамику решения задачи. Метод DyLAN~\cite{liu2023}
предлагает динамический выбор состава команды агентов на
каждом шаге, но варьирует только участников --- сама
структура связей между ними остается фиксированной. Метод
Mixture-of-Agents~\cite{moa2024} организует агентов
в многослойную ансамблевую конфигурацию с агрегацией ответов,
которая также не меняется во время решения.

Промежуточные между однозвенными и мульти-агентными решения
представлены работами Tree-of-Thoughts~\cite{tot2023} и
CAMEL~\cite{camel2023}. Первая разворачивает дерево вариантов
рассуждения с поиском в ширину, вторая моделирует ролевой
диалог двух агентов. В обоих случаях структура взаимодействия
определяется заранее и фиксируется на протяжении всей задачи.
Ближе к проблеме адаптации структуры подходят работы по дебатам
между языковыми моделями~\cite{du2023}, где несколько моделей
обмениваются аргументами для повышения качества вывода, однако
число оппонентов и протокол обмена в них зафиксированы заранее.

Общая характеристика рассмотренных работ заключается в том,
что выбор топологии в них производится либо архитектурно
(на этапе проектирования системы), либо заранее (на этапе
получения новой задачи). Возможность переключения топологии
в процессе решения, в зависимости от фазы и наблюдаемых
сигналов, в этих работах не рассматривается.

\subsection{Механизмы адаптации топологий}
\label{sec:rev:adaptation}

Адаптация мульти-агентных систем во время выполнения исследуется
существенно реже. В проанализированной выборке преобладают два
направления адаптации, ни одно из которых не покрывает
интересующий нас случай динамического переключения топологии между
фазами решения.

Первое направление --- адаптация состава команды. Уже упоминавшийся
метод DyLAN~\cite{liu2023} выбирает на каждом шаге подмножество
агентов, наиболее подходящее для текущего контекста. Аналогично
работа AgentVerse~\cite{agentverse2023} использует механизм
найма агентов, при котором система формирует команду под задачу
из заранее заданного пула. В обоих случаях изменяется не структура
взаимодействия, а множество участников --- сама схема обмена
сообщениями остается фиксированной.

Второе направление --- оптимизация конфигурации перед запуском
задачи. Системы GPTSwarm~\cite{zhuge2024} и
G-Designer~\cite{zhang2024} принимают решение о структуре графа
агентов до начала вычислений и не пересматривают его в дальнейшем.
Такие методы можно рассматривать как однократную офлайн-адаптацию,
которая полезна для классов однородных задач, но недостаточна
для задач со сложной фазовой структурой.

Сложные задачи неоднородны по своей внутренней структуре. Фаза
планирования может требовать дебатного формата для генерации
альтернатив, фаза исполнения --- иерархического разделения
труда, фаза проверки --- независимой рецензии. Отсюда следует
естественная гипотеза --- статическая топология, оптимальная
для усредненной картины, может быть субоптимальной для каждой
отдельной фазы. Эта гипотеза мотивирует разработку
механизма переключения топологии в процессе выполнения задачи,
который и является одним из центральных предметов настоящего
исследования.

\subsection{Безопасность и устойчивость топологий}
\label{sec:rev:safety}

Отдельное направление работ посвящено анализу уязвимостей
мульти-агентных систем и зависимости этих уязвимостей от
топологии. Работа NetSafe~\cite{yu2024} систематически исследует
устойчивость различных конфигураций к воздействию скомпрометированного
агента. Авторы устанавливают, что плотные топологии (полносвязная,
звезда) более подвержены каскадному распространению ошибок, тогда
как разреженные структуры (цепочка) ограничивают радиус
поражения. Работа TOMA~\cite{liang2024} развивает анализ на
случай многошаговых атак, в которых злоумышленник может
последовательно воздействовать на несколько агентов с учетом
их связности.

Эти результаты важны для нашего исследования по двум причинам.
Во-первых, они показывают, что выбор топологии не сводится
к чистому критерию качества решения --- надежность системы
к ошибкам отдельных компонентов также является функцией
структуры графа. Во-вторых, они мотивируют включение в систему
защитных правил переключения топологий, которые предотвращают
патологические режимы взаимодействия (в частности, частое
осцилляционное переключение между топологиями без существенного
прогресса в решении). Реализация таких правил в настоящей работе
обсуждается в главе~\ref{sec:architecture}.

\subsection{Человек в контуре управления}
\label{sec:rev:hitl}

Подключение человека к работе автоматизированных систем (от англ.
Human-in-the-Loop, HITL) --- давно установленная практика. Обзорная
работа~\cite{mosqueira2023} систематизирует методы интеграции
человека в контур машинного обучения, выделяя такие функции
человека как разметка данных, исправление предсказаний, активный
отбор примеров. Классическая работа Хорвица~\cite{horvitz1999}
ввела принципы смешанной инициативы, при которой автоматизированная
система и человек поочередно ведут процесс решения задачи в
зависимости от уверенности каждой стороны. Эти принципы остаются
актуальными для проектирования современных мульти-агентных
систем.

Отдельную линию исследований представляют методы обучения
языковых моделей с обратной связью от человека. Работа
Кристиано и соавторов~\cite{christiano2017} предложила обучение
с подкреплением от человеческих предпочтений, что позднее легло
в основу метода InstructGPT~\cite{ouyang2022} и стало стандартной
техникой выравнивания моделей. В этих работах человек выступает
в роли поставщика сигнала вознаграждения. Применительно к
мульти-агентным системам вопрос ставится принципиально иначе ---
человек участвует не только в обучении, но и в каждом отдельном
вызове системы, что требует формализации его роли в графе
обмена сообщениями.

Обзор~\cite{wang2024} перечисляет несколько типичных паттернов
включения человека --- координатор, рецензент, судья при споре
агентов, --- но не предлагает экспериментального сравнения их
эффективности. Фреймворк AutoGen~\cite{wu2023} допускает
подключение человека как агента в диалоге, однако роль человека
задается архитектурно и не варьируется в ходе экспериментов.
Отдельные работы исследуют человека как судью в задачах оценки
качества генерации с использованием шкалы NASA-TLX (от англ.
NASA Task Load Index, индекс рабочей нагрузки NASA)~\cite{nasatlx1988},
но эта функция выделена в самостоятельную ветвь литературы и
не пересекается с проектированием рабочего графа агентов.

Современные оркестраторы мульти-агентных систем, такие как
Magentic-One~\cite{magenticone2024}, реализуют разделение между
ведущим агентом-координатором и подчиненными работниками, при
этом человек явно подключен как опциональный участник в роли
постановщика задачи. Однако сравнения эффективности разных ролей
человека в этой работе также не проводится.

В проанализированной выборке работ за 2022--2026 годы
систематических исследований того, какая роль человека и в какой
фазе наиболее эффективна для разных типов задач, не обнаружено.
Эта лакуна составляет второй центральный предмет настоящего
исследования.

\subsection{Сравнительная аналитическая таблица методов}
\label{sec:rev:comparison}

Рассмотренные методы естественно сопоставить по нескольким
критериям, релевантным для исследовательских вопросов работы.
Соответствующее сравнение приведено в таблице~\ref{tab:rev:methods}.
Знак \enquote{$+$} означает наличие свойства, \enquote{$-$} ---
его отсутствие, \enquote{$\pm$} --- частичную реализацию.

\begin{table}[!htbp]
\footnotesize
\caption{Сравнение существующих методов по ключевым критериям}
\label{tab:rev:methods}
\footnotesize
\begin{tabularx}{\textwidth}{l c c c c c c c}
\toprule
\textbf{Метод} & \textbf{Год} & \textbf{MAS} & \textbf{Сравн.}
& \textbf{Адапт.} & \textbf{Адапт.} & \textbf{HITL} & \textbf{Безопасн.} \\
& & & \textbf{топ.} & \textbf{состава} & \textbf{топ. runtime} & & \\
\midrule
ReAct~\cite{react2022}             & 2022 & $-$ & $-$ & $-$ & $-$ & $-$ & $-$ \\
AgentVerse~\cite{agentverse2023}   & 2023 & $+$ & $-$ & $+$ & $-$ & $-$ & $-$ \\
LLM debate~\cite{du2023}           & 2023 & $+$ & $-$ & $-$ & $-$ & $-$ & $-$ \\
CAMEL~\cite{camel2023}             & 2023 & $+$ & $-$ & $-$ & $-$ & $-$ & $-$ \\
DyLAN~\cite{liu2023}               & 2023 & $+$ & $-$ & $+$ & $-$ & $-$ & $-$ \\
Reflexion~\cite{reflexion2023}     & 2023 & $-$ & $-$ & $-$ & $-$ & $-$ & $-$ \\
AutoGen~\cite{wu2023}              & 2023 & $+$ & $-$ & $-$ & $-$ & $+$ & $-$ \\
Tree-of-Thoughts~\cite{tot2023}    & 2023 & $-$ & $-$ & $-$ & $-$ & $-$ & $-$ \\
Magentic-One~\cite{magenticone2024}& 2024 & $+$ & $-$ & $-$ & $-$ & $\pm$ & $-$ \\
MetaGPT~\cite{hong2024}            & 2024 & $+$ & $-$ & $-$ & $-$ & $-$ & $-$ \\
TOMA~\cite{liang2024}              & 2024 & $+$ & $\pm$ & $-$ & $-$ & $-$ & $+$ \\
Self-Refine~\cite{selfrefine2024}  & 2024 & $-$ & $-$ & $-$ & $-$ & $-$ & $-$ \\
ChatDev~\cite{qian2024}            & 2024 & $+$ & $-$ & $-$ & $-$ & $-$ & $-$ \\
Scaling MAS~\cite{qian2024scaling} & 2024 & $+$ & $\pm$ & $-$ & $-$ & $-$ & $\pm$ \\
MoA~\cite{moa2024}                 & 2024 & $+$ & $-$ & $-$ & $-$ & $-$ & $-$ \\
NetSafe~\cite{yu2024}              & 2024 & $+$ & $+$ & $-$ & $-$ & $-$ & $+$ \\
G-Designer~\cite{zhang2024}        & 2024 & $+$ & $+$ & $-$ & $-$ & $-$ & $-$ \\
GPTSwarm~\cite{zhuge2024}          & 2024 & $+$ & $\pm$ & $-$ & $-$ & $-$ & $-$ \\
\midrule
\textbf{Настоящая работа}          & 2026 & $+$ & $+$ & $-$ & $+$ & $+$ & $\pm$ \\
\bottomrule
\end{tabularx}
\\[0.3cm]
\footnotesize
Обозначения колонок ---
\textbf{MAS} мульти-агентная архитектура,
\textbf{Сравн. топ.} систематическое сравнение нескольких топологий
на едином стенде,
\textbf{Адапт. состава} динамическое изменение участников,
\textbf{Адапт. топ. runtime} переключение структуры графа в процессе
решения задачи,
\textbf{HITL} интеграция человека в контур управления как
полноценного участника графа,
\textbf{Безопасн.} учет устойчивости к скомпрометированному агенту
и каскадным сбоям. Знак \enquote{$+$} означает наличие свойства,
\enquote{$-$} --- его отсутствие, \enquote{$\pm$} --- частичную
реализацию.
\end{table}

\subsection{Что отсутствует в существующих работах}
\label{sec:rev:gap}

Из таблицы~\ref{tab:rev:methods} видно, что три столбца ---
систематическое сравнение топологий на едином стенде, адаптация
топологии во время решения задачи, явная роль человека в
структуре графа --- в существующих работах представлены лишь
эпизодически и в проанализированной выборке не встречаются
одновременно. Метод G-Designer~\cite{zhang2024}
систематически сравнивает топологии, но не адаптирует их во
время выполнения задачи. AutoGen~\cite{wu2023} позволяет подключать
человека, но не формализует его роль и не сравнивает разные
конфигурации. DyLAN~\cite{liu2023} адаптирует состав, но не
структуру связей. Ни один из методов не объединяет в одном
эксперименте сравнение нескольких топологий, механизм их
переключения во время решения и многомерное варьирование роли
человека.

Эта совокупность отсутствующих свойств и составляет
исследовательскую лакуну, заполнение которой является целью
настоящей работы. В частности, настоящая работа предлагает ---
во-первых, единый программный стенд, на котором сравниваются
пять различных топологий на едином корпусе из четырех классов
задач; во-вторых, мета-топологию, переключающую активную базовую
топологию в зависимости от фазы решения; в-третьих,
формализацию пяти ролей человека и экспериментальную оценку их
влияния на каждую топологию по полному факторному плану.

\subsection{Выводы по главе}
\label{sec:rev:summary}

Проведенный аналитический обзор позволяет сделать несколько
содержательных выводов. Мульти-агентные системы больших языковых
моделей являются активно развивающейся областью с устоявшимся
набором архитектурных образцов, однако коммуникационная топология
в них обычно фиксируется на этапе проектирования и не пересматривается
во время выполнения задачи. Существующие механизмы адаптации
ограничены либо составом команды, либо однократным офлайн-выбором
структуры графа, что недостаточно для задач со сложной фазовой
структурой. Безопасность мульти-агентных систем существенно зависит
от топологии, что подтверждает необходимость защитных правил при
ее динамическом изменении. Систематического исследования роли
человека в~графе агентов в~проанализированной выборке работ
не обнаружено. Среди смежных работ по~бенчмаркингу агентских систем
(AgentBench~\cite{agentbench2023}) и~по~маршрутизации между
языковыми моделями (RouteLLM~\cite{routellm2024},
FrugalGPT~\cite{frugalgpt2023}) формализованной адаптивной роли
человека также не введено.

На пересечении трех направлений --- сравнение топологий, их
динамическая адаптация, формализованная роль человека --- лежит
лакуна, заполнение которой ставится в качестве научной цели
настоящей работы. Эта цель формализуется в виде четырех
исследовательских вопросов в разделе~\ref{sec:methodology}.


\newpage
% =====================================================================
% Глава 2. Постановка задачи и методология исследования.
% Объем 5--6 страниц.
% =====================================================================

\section{Постановка задачи и методология исследования}
\label{sec:methodology}

Глава формализует объект исследования, исследовательские вопросы
и методологию их эмпирической проверки --- математическая модель,
шесть топологий, пять ролей человека, система метрик, обоснование
корпуса и моделей, воронка из четырех экспериментов.

\subsection{Формальная модель мульти-агентной системы}
\label{sec:meth:formal}

Мульти-агентная система рассматривается как кортеж
\begin{equation}
\mathcal{M} = \langle V,\, E_t,\, R,\, \mathcal{T},\, F,\, h \rangle,
\end{equation}
где $V = \{v_1, \dots, v_n\}$ --- конечное множество агентов,
каждый из которых реализован как языковая модель с фиксированной
системной подсказкой и набором инструментов; $E_t \in 2^{V \times V}$
--- множество допустимых направленных ребер коммуникации на
итерации~$t$, образующее коммуникационный граф; $R \colon V \to
\mathcal{R}_A$ --- отображение агентов в множество ролей агентов
$\mathcal{R}_A$ из шести значений (планировщик, исследователь,
исполнитель, критик, дебатер, координатор); $\mathcal{T}$ ---
текущая задача с целевой функцией оценки качества решения;
$F$ --- конечное множество из четырех фаз решения задачи
(планирование, исполнение, проверка, завершение); $h$ ---
(опционально) участник-человек с~ролью из~отдельного множества
$\mathcal{R}_H$ (см.~\ref{sec:meth:roles}).

В классических работах граф $E_t$ определяется архитектурно и не
меняется в ходе выполнения задачи. В настоящем исследовании
вводится расширение модели --- структура $E_t$ становится результатом
оператора обновления, зависящего от~текущей фазы $f_t \in F$ и
наблюдаемых сигналов $\sigma_t \in \{0,1\}^m$ агентов:
\begin{equation}
E_t = \Phi(f_t,\, \sigma_t,\, t),
\label{eq:meth:adaptive}
\end{equation}
где $\Phi$ --- оператор обновления структуры коммуникационного
графа, $m$ --- фиксированная размерность вектора сигналов
(флаги завершения этапа, признаки тупикового состояния, счетчики
отказов критика и другие наблюдаемые индикаторы прогресса).
Возможность такого переключения и составляет суть исследуемой
адаптивной мета-топологии.

Множество ролей агентов $\mathcal{R}_A$ выше уже зафиксировано
шестью значениями. Множество ролей человека~$\mathcal{R}_H$
фиксировано и содержит пять значений, формализованных в
разделе~\ref{sec:meth:roles}. Фазы $F$ упорядочены отношением
$\textit{planning} \prec \textit{execution} \prec
\textit{verification} \prec \textit{done}$ с монотонной семантикой
переходов.

Задача оптимизации мульти-агентной системы формулируется как
многокритериальная по~двум первичным осям~--- качеству решения
$q(\mathcal{M})$ и~суммарной стоимости вызовов $c(\mathcal{M})$~---
при~ограничении на~время выполнения
$\tau(\mathcal{M}) \le \tau_{\max}$, где $\tau_{\max}$ задается
конфигурацией эксперимента. Поскольку чистый Парето-оптимум обычно
не единственен, в работе используется скаляризация с фиксированными
порогами улучшения по~двум первичным осям. Конфигурация
$\mathcal{M}_1$ считается выигрышной относительно $\mathcal{M}_0$,
если соблюдено ограничение $\tau(\mathcal{M}_1) \le \tau_{\max}$ и
выполняется одно из условий
\begin{equation}
\frac{q(\mathcal{M}_1) - q(\mathcal{M}_0)}{q(\mathcal{M}_0)}
   \geq 0{,}05
   \quad\text{или}\quad
\frac{c(\mathcal{M}_0) - c(\mathcal{M}_1)}{c(\mathcal{M}_0)}
   \geq 0{,}15
   \text{ при } q(\mathcal{M}_1) \geq q(\mathcal{M}_0).
\label{eq:meth:thresholds}
\end{equation}
Пороги $5\%$ для качества и $15\%$ для стоимости зафиксированы
до проведения основных запусков и обоснованы в
разделе~\ref{sec:meth:metrics}.

\subsection{Шесть рассматриваемых топологий}
\label{sec:meth:topologies}

Пространство сравнения составляют шесть топологий. Пять из них
являются статичными и реализуют пять базовых классов коммуникационных
структур, перечисленных в обзоре литературы. Шестая является
мета-топологией, переключающей активную статическую топологию
в зависимости от фазы.

\textbf{Star} --- центральный координатор $v_c$ последовательно
обращается к специализированным работникам $v_1, \dots, v_k$ и
агрегирует их ответы. Граф ребер $E_{\text{star}} =
\{(v_c, v_i), (v_i, v_c) \mid i = 1, \dots, k\}$.

\textbf{Chain} --- линейный конвейер из трех агентов
$v_p \to v_e \to v_c$ (планировщик, исполнитель, критик)
с условным циклом доработки. Граф направленный ациклический
с обратной связью от критика к исполнителю.

\textbf{Mesh} --- полносвязная топология с общей шиной сообщений.
Граф $E_{\text{mesh}} = (V \times V) \setminus \{(v,v)\}$, обмен
организован по принципу циклической очереди, решение принимается
голосованием.

\textbf{Debate} --- пара оппонентов $v_+$ и $v_-$ генерирует
параллельные альтернативные решения, после чего судья $v_j$
выбирает победителя или синтезирует ответ. Граф двудольный
с агрегатором.

\textbf{Hierarchical} --- двухуровневая иерархия из координатора
верхнего уровня $v_c^{\text{top}}$ и двух подкоманд
$T_a, T_b \subset V$, каждая из которых сама по себе является
графом. Подкоманды работают параллельно, координатор синтезирует
итог.

\textbf{Adaptive} --- мета-топология, активирующая на каждой
итерации одну из пяти базовых конфигураций. Решение о выборе
принимается маршрутизатором согласно
формуле~(\ref{eq:meth:adaptive}). Псевдокод процедуры приведен
в таблице~\ref{tab:meth:adaptive}, общая концептуальная схема
взаимодействия блоков системы изображена на
рисунке~\ref{fig:arch:functional} главы~\ref{sec:architecture},
архитектура реализации описана в разделе~\ref{sec:arch:topology}.

\begin{table}[!htbp]
\caption{Процедура работы адаптивной мета-топологии}
\label{tab:meth:adaptive}
\small
\begin{tabularx}{\textwidth}{r X}
\toprule
\textbf{Шаг} & \textbf{Действие} \\
\midrule
1  & Инициализировать фазу $f \leftarrow planning$, начальную
     топологию $K \leftarrow K_0$, состояние $s \leftarrow s_0$
     и журнал $\mathcal{J} \leftarrow \varnothing$. \\
2  & Пока $f \neq done$ и $c(s) < B$ выполнять шаги 3--10. \\
3  & Применить маршрутизатор фаз
     $f' \leftarrow \mathrm{PhaseRouter}(s, f)$. \\
4  & Применить маршрутизатор топологий
     $K' \leftarrow \mathrm{TopologyRouter}(s, K, f')$. \\
5  & Если $\mathrm{Guard}(s, K, K') = \mathit{block}$, то
     $K' \leftarrow K$ и записать $\mathit{guard\_override}$
     в $\mathcal{J}$. \\
6  & Если $K' \neq K$, применить ворота перехода состояния
     $s \leftarrow \mathrm{TransitionGate}(s, K, K')$ и записать
     переход в $\mathcal{J}$. \\
7  & Исполнить подграф выбранной топологии
     $s \leftarrow \mathrm{Subgraph}_{K'}(s)$. \\
8  & Обновить $K \leftarrow K'$ и $f \leftarrow f'$. \\
9  & Записать журнал итерации в $\mathcal{J}$. \\
10 & Перейти к шагу 2. \\
11 & Вернуть результат $y(s)$ и журнал $\mathcal{J}$. \\
\bottomrule
\end{tabularx}
\end{table}

\subsection{Пять ролей человека в контуре управления}
\label{sec:meth:roles}

Пять ролей человека формализованы по аналогии с типовыми позициями
в человеческих организациях. \textbf{Координатор} ($\rho_C$) ---
управляет планом и делегированием, активен в иерархических и
полносвязных топологиях. \textbf{Рецензент} ($\rho_R$) --- проверяет
промежуточные результаты и инициирует доработку, естественно
вписывается в Star и Chain. \textbf{Судья} ($\rho_J$) ---
выносит финальный вердикт при наличии альтернатив, релевантен
для Debate. \textbf{Равноправный участник} ($\rho_P$) ---
вносит альтернативные мнения наравне с агентами, вписывается
в Mesh. \textbf{Наблюдатель} ($\rho_M$) --- следит за процессом
и вмешивается только в критических случаях (превышение бюджета,
зависание, циклы), применим в любой топологии.

Не все пары (топология, роль) являются осмысленными.
Рассматриваются пять статических топологий и пять ролей человека ---
полное произведение содержит 25 пар. Из него исключаются три
заведомо вырожденные комбинации. Пара (Debate, Coordinator)
исключается, поскольку координатор в дебатной топологии
эквивалентен судье. Пары (Chain, Peer) и (Hierarchical, Peer)
исключаются, поскольку равноправный участник ломает линейную или
иерархическую структуру. После исключений пространство содержит
$5 \times 5 - 3 = 22$ осмысленные комбинации
$\text{топология} \times \text{роль человека}$, где~используются
пять статических топологий и~пять ролей человека из~$\mathcal{R}_H$.

\subsection{Система оценочных метрик}
\label{sec:meth:metrics}

Метрики разделены на три группы. Первая группа --- базовые метрики
качества и эффективности --- применима ко всем экспериментам и
ко всем топологиям.

\begin{itemize}[noitemsep]
  \item $q \in [0, 1]$ --- агрегированный показатель качества
        решения. Конкретная формула зависит от класса задач ---
        для задач программирования это метрика pass@1 (доля
        задач, для которых сгенерированный код проходит все
        предзаданные модульные тесты с первой попытки); для
        математических задач --- exact match (точное совпадение
        числового ответа с эталоном); для творческой генерации ---
        полусумма метрики ROUGE-L~\cite{rouge2004} (от англ.
        Recall-Oriented Understudy for Gisting Evaluation, Longest
        common subsequence variant) и~доли покрытых концептов; для
        анализа данных --- exact match числового ответа.
  \item $c$ --- суммарная стоимость вызовов языковых моделей
        в долларах США за один запуск.
  \item $\tau$ --- время выполнения запуска по настенным часам
        (от англ. wall-clock) в секундах.
  \item $r_{\text{fail}}$ --- доля запусков, завершившихся со
        статусом, отличным от успешно выполненного.
  \item $N_{\text{iter}}$ --- число итераций основного цикла.
\end{itemize}

Вторая группа --- метрики, специфичные для адаптивной мета-топологии.

\begin{itemize}[noitemsep]
  \item $N_{\text{switch}}$ --- число фактических переключений
        активной топологии за один запуск.
  \item $r_{\text{guard}}$ --- доля решений маршрутизатора,
        заблокированных защитными правилами.
  \item $\gamma$ --- доля бюджета, потраченная маршрутизатором
        на собственные вызовы языковой модели; критичная метрика,
        отвечающая на вопрос, не превышает ли стоимость адаптации
        получаемый от нее выигрыш.
  \item $r_{\text{hurt}}$ --- доля задач, на которых адаптивная
        конфигурация дает качество хуже лучшей статической
        конфигурации. Является ключевым индикатором безопасности
        механизма адаптации.
  \item $\Delta_{\text{oracle}}$ --- разрыв качества между
        теоретическим per-task best-of референсом и~фактическим
        маршрутизатором, формально определяемый
        по~формуле~(\ref{eq:meth:delta-oracle}). Per-task best-of
        референс вычисляется по~результатам базового эксперимента~E1~---
        для~каждого корпуса задач выбирается статическая топология
        с~максимальным средним качеством на~этом корпусе. Референс
        служит верхней границей достижимого качества для~класса
        задач при~заданном пространстве статических топологий и~требует
        апостериорного знания пары $(\text{корпус},\text{лучшая топология})$;
        в~условиях открытого развертывания без~классификатора задач недостижим.
\end{itemize}

\begin{equation}
\Delta_{\text{oracle}} = q_{\text{best-of}} - q_{\text{router}},
\quad\text{где}\quad
q_{\text{best-of}} = \frac{1}{|\mathcal{C}|}
   \sum_{c \in \mathcal{C}}
      q\bigl(K^*_c,\, c \bigr),
\label{eq:meth:delta-oracle}
\end{equation}

\noindent где $\mathcal{C}$ --- множество корпусов задач,
$K^*_c = \arg\max_K q(K, c)$ --- оптимальная статическая топология
для~корпуса~$c$, $q(K, c)$ --- среднее качество топологии~$K$
на~корпусе~$c$ по~базовому эксперименту~E1.

Третья группа --- метрики, относящиеся к взаимодействию с человеком.

\begin{itemize}[noitemsep]
  \item $N_{\text{int}}$ --- число обращений системы к человеку
        за один запуск.
  \item $L_{\text{cog}}$ --- прокси-метрика когнитивной нагрузки
        для запусков с симулятором человека, формально определяемая
        по формуле~(\ref{eq:meth:lcog}). В экспериментах с реальными
        участниками заменяется на стандартную шкалу NASA-TLX.
\end{itemize}

\begin{equation}
L_{\text{cog}} = w_n \cdot \widetilde{N}_{\text{int}}
              + w_l \cdot \widetilde{L}_{\text{ctx}}
              + w_\tau \cdot \widetilde{\tau}_{\text{resp}},
\quad
w_n = w_l = w_\tau = \tfrac{1}{3},
\label{eq:meth:lcog}
\end{equation}

\noindent где $\widetilde{N}_{\text{int}}$, $\widetilde{L}_{\text{ctx}}$
и $\widetilde{\tau}_{\text{resp}}$ --- нормированные на отрезок
$[0,1]$ значения числа обращений к участнику, средней длины
предоставляемого ему контекста и средней задержки ответа симулятора
соответственно. Нормировка выполняется по квантилям распределения
наблюдаемых значений за все запуски эксперимента. Равные веса
выбраны априорно до начала запусков как нейтральная свертка трех
независимых компонентов нагрузки.

Адаптивная конфигурация считается выигрышной относительно
наилучшей статической, если она дает прирост качества не менее
$5\%$ либо снижение стоимости не менее $15\%$ при сохранении
качества. Пороги фиксируются предварительно, до проведения
основных запусков, что соответствует практике предварительной
регистрации гипотез в эмпирических исследованиях.

\subsection{Корпус задач и обоснование его выбора}
\label{sec:meth:corpus}

Эксперименты проводятся на стратифицированной выборке из четырех
открытых корпусов, покрывающих четыре качественно различных
класса задач --- программирование, математические рассуждения,
творческая генерация и анализ данных. Сравнение корпуса
с альтернативами приведено в таблице~\ref{tab:meth:corpora}.

\begin{table}[!htbp]
\caption{Сравнение рассмотренных корпусов задач}
\label{tab:meth:corpora}
\footnotesize
\renewcommand{\arraystretch}{1.25}
\begin{tabularx}{\textwidth}{%
   >{\raggedright\arraybackslash}p{2.6cm}%
   >{\raggedright\arraybackslash}p{2.4cm}%
   c%
   >{\raggedright\arraybackslash}p{1.7cm}%
   >{\raggedright\arraybackslash}X}
\toprule
\textbf{Корпус} & \textbf{Класс} & \textbf{Размер} &
\textbf{Метрика} & \textbf{Решение по~включению} \\
\midrule
HumanEval~\cite{humaneval2021}        & Программирование    & 164   & pass@1        & Включен; стандарт сравнения. \\
HumanEval+~\cite{humanevalplus2023}   & Программирование    & 164   & pass@1+       & Не~включен; превышает покрытие основной выборки. \\
LiveCodeBench~\cite{livecodebench2024} & Программирование   & 500+  & pass@1        & Резерв для~проверки на~загрязнение. \\
GSM8K~\cite{gsm8k2021}                & Мат. рассуждения    & 8500  & exact match   & Включен; стандарт оценки. \\
MATH~\cite{mathbench2021}             & Мат. рассуждения    & 12500 & exact match   & Не~включен; сложность выше возможностей рабочих моделей. \\
MMLU-Pro~\cite{mmlupro2024}           & Множ. рассуждения   & 12000 & accuracy      & Не~включен; формат ответа не~подходит для~дебатной топологии. \\
CommonGen~\cite{commongen2020}        & Творческая генерация & 35000 & ROUGE-L $+$ cov & Включен; покрывает свободная (без~однозначного ответа) генерацию с~ограничениями. \\
InfiAgent-DABench~\cite{dabench2024}  & Анализ данных       & 257   & numeric exact & Включен; единственный корпус с~замкнутым ответом по~аналитике данных. Качество чувствительный к~модели (см.~\ref{sec:exp:conf}). \\
HellaSwag~\cite{hellaswag2019}        & Здравый смысл       & 10000 & accuracy      & Не~включен; качество современных моделей близко к~потолку. \\
\bottomrule
\end{tabularx}
\end{table}

Выбор обусловлен совокупностью требований. Каждый класс задач
представлен ровно одним корпусом, чтобы избежать смешения
эффектов корпуса и класса. Каждый корпус допускает автоматическую
объективную оценку, исключающую субъективность оценщика
(юнит-тесты, точное совпадение числа, метрики покрытия концептов).
Размер задач в каждом корпусе позволяет уложиться в один запуск
рабочих моделей. От рассматриваемых для каждого класса альтернатив
выбранные корпуса отличаются распространенностью в литературе и
зрелостью методики оценки, что обеспечивает сопоставимость
результатов с предшествующими работами. Из каждого корпуса
производится стратифицированная выборка по 15 задач, итого
60 задач в одном базовом запуске. Каждая задача повторяется
с тремя различными зернами генерации (для контроля стохастичности
выборки токенов), что в комбинации с пятью статическими топологиями
дает $60 \times 3 \times 5 = 900$ запусков в эксперименте E1.

\subsection{Языковые модели и инфраструктура}
\label{sec:meth:models}

Распределение моделей по ролям изолирует переменную «модель» от
переменных «топология» и «роль человека». Все рабочие агенты
(планировщик, исследователь, исполнитель, критик, дебатер), а также
маршрутизаторы и симулятор человека используют единую модель
Cerebras gpt-oss-120b (~3000 токенов/с на Cerebras WSE). Иерархия
«работник слабее критика» реализуется различиями в подсказках и
параметре строгости рассуждения, а не разными моделями ---
фиксированная модельная мощность критична для чистоты сравнения
топологий.

Судья оценки качества вынесен в отдельную линию --- OpenAI
gpt-4.1-mini с самосогласованностью из трех вызовов и нулевой
температурой; модель из другого семейства весов исключает петлю
положительной обратной связи. Подтверждающие запуски выполняются
на Cerebras qwen-3-235b как представителе другого семейства весов.

Выбор моделей опирался на четыре критерия --- отказоустойчивый
программный интерфейс (от англ. Application Programming Interface,
API) с вызовом инструментов и структурированным выводом, низкая
удельная стоимость вызова, высокая пропускная способность и
доступность из РФ.

Инфраструктура (глава~\ref{sec:architecture}) --- LangGraph,
PostgreSQL~16, Apache Parquet; параллелизм --- 12 процессов с
асинхронной обработкой внутри каждого.

\subsection{Дизайн экспериментов --- воронка}
\label{sec:meth:funnel}

Прямой эксперимент по всему декартову произведению (6 топологий
$\times$ 5 ролей $\times$ 3 режима $\times$ 4 класса задач) дал бы
более 10\,000 ячеек. Вместо этого исследование организовано как
воронка из четырех экспериментов, каждый из которых сужает
пространство конфигураций по результатам предыдущего
(таблица~\ref{tab:meth:funnel}).

\begin{table}[!htbp]
\caption{Структура экспериментальной воронки}
\label{tab:meth:funnel}
\small
\renewcommand{\arraystretch}{1.25}
\begin{tabularx}{\textwidth}{%
   l%
   c%
   >{\raggedright\arraybackslash}p{3.0cm}%
   >{\raggedright\arraybackslash}X}
\toprule
\textbf{Эксп.} & \textbf{Вопрос} & \textbf{Запуски} &
\textbf{Цель} \\
\midrule
E1 & RQ1 & $5\cdot4\cdot15\cdot3=900$ & Парето-фронт пяти статических
   топологий на~четырех классах задач. \\
\addlinespace[3pt]
E2 & RQ3 & $\approx 2295$ (после фильтра top-3) & Влияние пяти ролей
   человека на~каждую топологию из~top-3~E1. \\
\addlinespace[3pt]
E3 & RQ2 & $\approx 640$ (237 адаптивных $+$ 406 статических) &
   Адаптивная мета-топология в~двух режимах маршрутизации
   (\code{rule}, \code{llm}) и~три статические базовые линии. \\
\addlinespace[3pt]
E4 & RQ4 & $3\cdot4\cdot15\cdot3=540$ & Адаптивная роль человека
   в~сочетании с~адаптивной топологией. \\
\addlinespace[3pt]
Conf. & --- & $120+142\approx 260$ & Кросс-семейная валидация
   выводов E3 и~E4 на~подтверждающей модели Qwen-3-235B. \\
\bottomrule
\end{tabularx}
\end{table}

Исследовательские вопросы формулируются следующим образом.

\textbf{RQ1.} Различаются ли пять статических топологий
по~количественным показателям качества, стоимости и~времени
для~разных классов задач?

\textbf{RQ2.} Раскладывается на~две взаимодополняющих части. (а)
Дает ли адаптивная мета-топология значимый прирост качества или~экономию
стоимости по~сравнению с~per-task best-of референсом и~с~отдельными
статическими топологиями? (б) Устойчив ли выбор оптимальной
стратегии маршрутизации (правиловой или~языковой) к~смене семейства весов
рабочего агента, или~он является решение, принимаемое для~каждого семейства моделей?

\textbf{RQ3.} Какая из~пяти ролей человека наиболее эффективна для
каждой топологии и~для~каждого класса задач?

\textbf{RQ4.} Превосходит ли совместная адаптация топологии и~роли
человека по~фазам решения фиксированную роль; сохраняется ли
направление эффекта при~смене семейства весов рабочего агента?

Статистическая обработка включает двухфакторный дисперсионный
анализ (two-way ANOVA, тип~III сумм квадратов для несбалансированного
дизайна) с~эффектом топологии и~эффектом роли в~качестве факторов
и~их~взаимодействием на~уровне значимости $\alpha = 0{,}05$.
При~нарушении допущений нормальности или равенства дисперсий
применяется непараметрический аналог по~критерию Краскела--Уоллиса.
Парные сравнения средних выполняются с~поправкой
Бенджамини--Хохберга~\cite{benjamini1995} на~множественное сравнение
при~уровне ложных открытий $\text{FDR} = 0{,}05$. Оценка практической
значимости различий проводится через коэффициент размера эффекта
Коэна~$d$~\cite{cohen1988}. Доверительные интервалы (от англ.
Confidence Interval, CI) рассчитываются по~методу непараметрического
бутстрапа~\cite{efron1979} (метод процентилей) с~десятью тысячами
ресэмплов на~уровне 95\,\% при~фиксированном начальном значении
генератора псевдослучайных чисел.

\subsection{Разграничение заимствованных и авторских элементов}
\label{sec:meth:boundary}

Граница между заимствованными и авторскими элементами. К
заимствованным относятся пять статических топологий, четыре
открытых корпуса задач, языковые модели и стандартный
статистический инструментарий (ANOVA, бутстрап, коэффициент Коэна).
К авторским --- адаптивная мета-топология вместе с процедурой
работы (\ref{eq:meth:adaptive}); формализация пяти ролей человека
и матрица их совместимости с топологиями; набор адаптивно-специфичных
метрик ($N_{\text{switch}}$, $r_{\text{guard}}$, $\gamma$,
$r_{\text{hurt}}$, $\Delta_{\text{oracle}}$); вороночная структура
исследования с предварительной фиксацией порогов
(\ref{eq:meth:thresholds}) и кросс-семейным подтверждением.

\subsection{Выводы и результаты по главе}
\label{sec:meth:summary}

Формализована MAS как кортеж из шести компонентов с динамическим
переключением графа по фазе; определены шесть топологий, пять
ролей человека и три группы метрик (включая адаптивно-специфичные).
Обоснован выбор корпуса из четырех бенчмарков и линии моделей с
кросс-семейным судьей. Сформулированы четыре RQ и воронка из
четырех экспериментов с подтверждающим запуском на дополнительном
семействе моделей. Граница заимствованного и авторского
зафиксирована в~\ref{sec:meth:boundary} и развернута
в~\ref{sec:contribution}. План реализуется в инфраструктуре
главы~\ref{sec:architecture} и проверяется в~\ref{sec:experiments}.


\newpage
% =====================================================================
% Глава 3. Архитектура программной системы.
% Версия 2 --- после применения правок независимого ревьюера.
% Объем 5--6 страниц.
% =====================================================================

\section{Архитектура программной системы}
\label{sec:architecture}

Программная инфраструктура работы реализована как фреймворк
адаптивных топологий для мульти-агентных систем
(от англ. Adaptive Topologies for Multi-agent systems, ATM)
на языке Python не ниже 3.11, распространяемый как пакет \code{atm}.
Глава описывает архитектуру с акцентом на решения, относящиеся к
научным гипотезам --- адаптивному переключению топологий и
интеграции человека. Исходный код --- в приложении~\ref{app:github}.

Функциональная схема --- на~рисунке~\ref{fig:arch:functional}. Вход ---
описание задачи и YAML-конфиг
(от англ. YAML Ain't Markup Language); выход --- журналы
взаимодействия, метрики и записи о фазах, переключениях топологий
и взаимодействиях с человеком. Между входом и выходом работают
пять блоков --- оркестратор, менеджер фаз и маршрутизаторы, активная
топология с агентами, шлюз HITL и подсистема наблюдаемости.

\begin{figure}[!htbp]
\centering
\fbox{\parbox{0.93\textwidth}{\centering\small
\vspace{0.3cm}
Вход\ \ $\longrightarrow$\ \ Оркестратор эксперимента\ \ $\longrightarrow$\\[0.1cm]
$\big[$\,Менеджер фаз\ \ $+$\ \ Маршрутизатор топологий\,$\big]$\\[0.1cm]
$\longrightarrow$\ \ Активная топология (Star / Chain / Mesh /
Debate / Hierarchical)\\[0.1cm]
$\longleftrightarrow$\ \ Агенты\ \ $\longleftrightarrow$\ \ Инструменты
\& песочница\\[0.1cm]
$\longleftrightarrow$\ \ Шлюз взаимодействия с человеком
(симулятор / интерактивный)\\[0.1cm]
$\longrightarrow$\ \ Подсистема наблюдаемости\ \ $\longrightarrow$\\[0.1cm]
Выход --- метрики, журналы, переходы.
\vspace{0.3cm}
}}
\caption{Функциональная схема предлагаемой мульти-агентной системы.
Прямоугольниками выделены архитектурные блоки фреймворка,
двунаправленные стрелки изображают многократные обмены сообщениями
в пределах одной итерации. Развернутая схема с указанием слоев
агентов, внешних служб и подсистемы хранения приведена в
приложении~\ref{app:uml}.}
\label{fig:arch:functional}
\end{figure}

\subsection{Обзор архитектуры и слои}
\label{sec:arch:overview}

Кодовая база организована в виде тринадцати функционально
изолированных слоев (пакетов высокого уровня), каждый из которых
отвечает за один аспект системы и содержит несколько программных
модулей. Перечень слоев и их назначений приведен в
таблице~\ref{tab:arch:layers}. Такая декомпозиция позволяет
независимо тестировать компоненты, заменять реализации (например,
провайдера языковой модели или шлюз человека в контуре управления
от англ. Human-in-the-Loop, HITL) и добавлять новые топологии без
изменения существующего кода.

\begin{table}[!htbp]
\caption{Функциональные слои фреймворка \code{atm}}
\label{tab:arch:layers}
\small
\begin{tabularx}{\textwidth}{l X}
\toprule
\textbf{Слой} & \textbf{Назначение} \\
\midrule
\code{core}          & Базовые типы и состояние графа --- сообщения, фазы,
                       роли, контейнеры состояния \code{AgentState},
                       \code{SharedState}, \code{GraphState}. \\
\code{llm}           & Унифицированная оболочка над провайдерами больших
                       языковых моделей, подсчет стоимости, бюджетные
                       ограничители, детерминированный \code{FakeLLM}
                       для тестов. \\
\code{tools}         & Реестр инструментов и песочница на~основе
                       контейнерной платформы Docker для
                       исполнения генерируемого кода. \\
\code{agents}        & Реализация пяти ролей агентов и базовый класс
                       \code{Agent} с управляющим циклом вызова
                       инструментов. \\
\code{topology}      & Шесть реализаций топологий поверх LangGraph
                       и общий протокол \code{Topology}. \\
\code{phases}        & Конечный автомат фаз, маршрутизатор топологий
                       и защитные правила переключения. \\
\code{human}         & Протокол \code{HumanGateway}, симулированный
                       и интерактивный шлюзы, маршрутизатор ролей
                       человека. \\
\code{storage}       & Объектно-реляционное отображение (от англ.
                       Object-Relational Mapping, ORM), асинхронные
                       сессии PostgreSQL, запись в формат Apache Parquet. \\
\code{observability} & Перехват событий LangGraph, сериализация
                       и отправка их в хранилище. \\
\code{tasks}         & Реализация и загрузчики бенчмарков
                       HumanEval, GSM8K, CommonGen, InfiAgent-DABench. \\
\code{evaluation}    & Судьи на основе языковой модели, оракулы
                       достоверности, агрегатор шкалы NASA-TLX
                       (от англ. NASA Task Load Index --- индекс
                       когнитивной нагрузки NASA). \\
\code{experiment}    & Запуск одиночных и грид-экспериментов,
                       схемы конфигов, командная строка. \\
\code{analysis}      & Загрузчики результатов, агрегаторы метрик,
                       построение графиков. \\
\bottomrule
\end{tabularx}
\end{table}

Архитектурный стиль опирается на три паттерна --- расширяемые точки
как протоколы Python (\code{Topology}, \code{HumanGateway},
\code{PhaseRouter}, \code{TopologyRouter}, \code{RoleRouter},
\code{Tool}) с \code{@runtime_checkable}; редукторное слияние
состояния графа при параллельном выполнении агентов с
идемпотентностью по идентификатору сообщения; реестры с декораторной
семантикой для топологий, инструментов и оракулов.

Технологический стек --- LangGraph~\cite{langgraph2024} поверх
LangChain~\cite{langchain2023} с~чекпойнтером для возобновления;
PostgreSQL~16 через асинхронные сессии SQL\-Alchemy~2.x для
метаданных и Apache Parquet для объемных журналов; YAML-конфиги
с~валидацией Pydantic~\cite{pydantic2024} поверх OmegaConf;
трассировка через LangSmith~\cite{langsmith2024}; линтер \code{ruff}
и анализатор типов \code{mypy} в строгом режиме.

Часть архитектурных решений заимствована из существующих работ
и фреймворков, часть является авторской разработкой. Точное
разграничение приведено в таблице~\ref{tab:arch:contribution}.

\begin{table}[!htbp]
\caption{Разделение архитектурных компонентов на заимствованные
и авторские}
\label{tab:arch:contribution}
\small
\begin{tabularx}{\textwidth}{>{\raggedright\arraybackslash}p{4.0cm}
                              >{\raggedright\arraybackslash}p{3.8cm}
                              >{\raggedright\arraybackslash}X}
\toprule
\textbf{Компонент} & \textbf{Источник} & \textbf{Авторский вклад} \\
\midrule
Оркестрация графов        & LangGraph (внешний) & Адаптация под мета-граф \\
Пять статичных топологий  & Прототипы из обзора литературы &
                                                     Унифицированная
                                                     реализация поверх
                                                     общего протокола \\
Адаптивная мета-топология & Своя разработка    & Полностью \\
Маршрутизаторы фаз        & Своя разработка    & Полностью, три режима
                                                 (rule-based, на основе
                                                 языковой модели, оракул) \\
Защитные правила переключения & Своя разработка & Полностью, четыре типа
                                                 ограничителей \\
Шлюз HITL                 & Своя разработка    & Полностью, три ролевых
                                                 маршрутизатора и политика
                                                 тайм-аутов \\
Оракул per-task best-of   & Своя разработка    & Полностью, в том числе
                                                 процедура построения
                                                 идеализированного референса \\
Подсистема наблюдаемости  & LangGraph callbacks & Адаптеры для PostgreSQL
                                                 и Apache Parquet, схема
                                                 событий \\
Бенчмарки                 & HumanEval, GSM8K, CommonGen, DABench &
                                                 Унифицированные загрузчики
                                                 и оценщики \\
Реализация шкалы NASA-TLX & Стандартная форма  & Прокси-агрегатор для
                                                 симулятора человека \\
\bottomrule
\end{tabularx}
\end{table}

\subsection{Слой агентов и оболочка над языковой моделью}
\label{sec:arch:agents}

Класс \code{Agent} в \code{atm.agents} реализует управляющий цикл ---
чтение входящих сообщений, формирование подсказки, вызов модели,
обработка инструментов и публикация ответа. Определены пять ролей
(\code{Planner}, \code{Researcher}, \code{Executor}, \code{Critic},
\code{Debater}) и дополнительный \code{Coordinator} для
иерархических и адаптивных топологий. Каждая роль задается
YAML-конфигом (подсказка, инструменты, окно сообщений, температура,
управление контекстом). Скрэтчпад --- скользящее окно с
автоматическим сжатием старых сообщений вспомогательной моделью;
политика выбрана по пилотному эксперименту.

Оболочка \code{LLMWrapper} в \code{atm.llm} унифицирует провайдеров
(OpenAI, Anthropic, Cerebras, локальный \code{FakeLLM}),
инкапсулирует подсчет токенов и стоимости (таблица \code{Pricing}),
экспоненциальный повтор и фиксацию версии модели. Ограничители
\code{BudgetGuard} работают на трех уровнях (вызов, запуск,
грид-эксперимент). \code{FakeLLM} поддерживает три режима ---
сценарный, эхо и воспроизведение по записи --- для детерминированного
повторения экспериментов.

\subsection{Реализация топологий поверх LangGraph}
\label{sec:arch:topology}

Каждая из шести топологий реализована классом, удовлетворяющим
протоколу \code{Topology} с фабричным методом \code{build},
возвращающим скомпилированный граф LangGraph. Граф связывает
узлы --- агентов и управляющие функции --- ребрами трех типов
(безусловные, условные, параллельные через примитив \code{Send}).
Все графы подключены к чекпойнтеру и перехватчику событий
\code{atm.observability}.

Пять статических топологий описаны в~\ref{sec:meth:topologies}.
Шестая, \emph{Adaptive}, --- мета-граф второго уровня. На каждом
тике маршрутизаторы фаз и топологий выбирают базовую топологию,
вызываемую как подграф; ворота \code{TransitionGate} применяют
правила передачи состояния и записывают \code{TopologyTransition}
в журнал. Подграфы кэшируются по имени для исключения повторной
компиляции; решения маршрутизаторов хранятся в замыкании, а не
в общем состоянии, чтобы избежать перезаписи во вложенном подграфе.

\subsection{Менеджер фаз и маршрутизатор топологий}
\label{sec:arch:phases}

Решение моделируется конечным автоматом (от англ. Finite-State
Machine, FSM) из четырех состояний (планирование, исполнение,
проверка, завершение). Маршрутизатор фаз в \code{atm.phases}
реализован в двух режимах --- \code{RuleBasedPhaseRouter} (таблица
сигналов с монотонной проверкой, фаза не движется назад) и
\code{LLMPhaseRouter} (парсинг ответа модели в формате
JSON --- от англ. JavaScript Object Notation --- с откатом на
правило при любой ошибке валидации).

Маршрутизатор топологий реализован в трех режимах. \code{RuleBased}
использует таблицу (фаза $\times$ сигналы)~$\to$~топология
(\code{stuck} в исполнении $\to$ Mesh, \code{rejected_count}~$\ge 3$
$\to$ Debate); пороги подобраны в пилотном эксперименте.
\code{LLMTopologyRouter} принимает решение по подсказке с описанием
состояния, с rule-based откатом. \code{OracleTopologyRouter} читает
таблицу решений из файла --- реализует идеализированный референс;
применяется только для построения идеализированного per-task best-of
референса и~в~экспериментальной воронке главы~\ref{sec:experiments}
как~самостоятельный режим маршрутизации не~запускается. Сравнение
RQ2 ведется между двумя реализуемыми в~условиях открытого развертывания
режимами~--- \code{RuleBased} и~\code{LLMTopologyRouter}.

Защитные правила (\code{SwitchGuards}) предотвращают
осцилляционный режим --- четыре ограничения: \code{min_dwell},
\code{cooldown}, лимиты за фазу и за запуск. Заблокированные
решения фиксируются как \code{guard_override} для последующей
оценки доли блокировок как метрики стабильности.

\subsection{Шлюз взаимодействия с человеком}
\label{sec:arch:human}

Интеграция человека абстрагирована протоколом \code{HumanGateway}
с асинхронным методом запроса ответа (идентификаторы запуска и
запроса, роль, история, допустимые действия). Идемпотентность по
паре (\code{run_id}, \code{request_id}) --- повторный вызов
возвращает ранее полученный ответ, что критично для возобновления
прерванных запусков.

Реализованы два шлюза --- \code{LLMSimulatedGateway} (языковая
модель с подсказкой по роли) и \code{CLIGateway} (интерактивный
интерфейс командной строки от англ. Command-Line Interface, CLI
для реальных участников). Маршрутизатор ролей выбирает
активную роль динамически (rule-based таблица или решение модели)
либо фиксированно (\code{FixedRoleRouter}). Политика тайм-аутов
имеет три варианта (откат к модели, пропуск, повтор).

\subsection{Хранение, командная строка и воспроизводимость}
\label{sec:arch:storage}

Метаданные --- в семи таблицах PostgreSQL (\code{experiments},
\code{runs}, \code{phases}, \code{topology_transitions},
\code{human_interactions}, \code{budget_events},
\code{llm_calls}) с индексами по типичным аналитическим запросам
и составным индексом (эксперимент, статус) для выборки
невыполненных запусков. Объемные журналы --- в файлах Apache
Parquet с асинхронным писателем и партиционированием по запуску.
Миграции --- Alembic.

CLI \code{atm} на библиотеке Typer объединяет семь команд
(таблица~\ref{tab:arch:cli}).

\begin{table}[!htbp]
\caption{Команды интерфейса командной строки фреймворка}
\label{tab:arch:cli}
\small
\begin{tabularx}{\textwidth}{l X}
\toprule
\textbf{Команда} & \textbf{Назначение} \\
\midrule
\code{atm run}       & Запуск одиночного эксперимента из YAML-конфига. \\
\code{atm grid}      & Параллельный грид-эксперимент через
                       \code{ProcessPoolExecutor}. \\
\code{atm estimate}  & Предварительная оценка стоимости запуска
                       по таблице цен и историческим данным. \\
\code{atm status}    & Агрегированный статус эксперимента ---
                       число выполненных, неудачных и активных
                       ячеек, суммарная стоимость. \\
\code{atm resume}    & Возобновление прерванного запуска из
                       последнего чекпойнта LangGraph. \\
\code{atm replay}    & Детерминированное воспроизведение запуска
                       через \code{FakeLLM} в режиме воспроизведения
                       по записи. \\
\code{atm reconcile} & Поиск и пометка запусков-зомби (запись о статусе
                       \enquote{выполняется}, но процесс уже
                       завершен). \\
\bottomrule
\end{tabularx}
\end{table}

Воспроизводимость --- инварианты, фиксируемые при старте запуска
(\code{seed_all} для Python/NumPy/провайдеров/генератора
идентификаторов; \code{git_sha} из \code{git rev-parse HEAD};
снимок версий моделей в \code{model_version_snapshot} по первому
ответу провайдера; дайджест образа Docker-песочницы). Совокупность
полей делает возможным сравнение запусков между собой и анализ
артефактов депрекации моделей.

\subsection{Выводы и результаты по главе}
\label{sec:arch:summary}

Инфраструктура реализована как фреймворк из тринадцати слоев на
Python с расширяемостью через протоколы, воспроизводимостью через
фиксацию инвариантов и наблюдаемостью через единый перехват
событий LangGraph; в нее входят шесть топологий, два режима фаз,
три режима топологий и два шлюза HITL --- все необходимое для
экспериментов главы~\ref{sec:experiments}. Разграничение
заимствованных и авторских компонентов --- в
таблице~\ref{tab:arch:contribution}, авторский вклад
развернут в~\ref{sec:contribution}. Кодовая база основного пакета
\code{atm} --- 21\,327 строк Python в 112 модулях внутри
тринадцати слоев; конфигурация --- 43 YAML-файла (24 описывают
запуски E1--E4).


\newpage
% =====================================================================
% Глава 4. Экспериментальное исследование.
% Версия 2 --- после получения данных E1--E4 + confirmation runs.
% Объем 8--10 страниц.
% =====================================================================

\section{Экспериментальное исследование}
\label{sec:experiments}

Настоящая глава содержит результаты эмпирической проверки
исследовательских вопросов (от англ. Research Questions, RQ)
RQ1--RQ4, поставленных в разделе~\ref{sec:meth:funnel}. Глава
организована как последовательное изложение четырех основных
экспериментов воронки E1--E4, двух подтверждающих прогонов
на~кросс-семейном рабочем агенте и сводного статистического анализа
на~основе непараметрического бутстрапа.

\subsection{Условия проведения и общие характеристики}
\label{sec:exp:setup}

Все эксперименты выполнены в единой программной инфраструктуре,
описанной в главе~\ref{sec:architecture}, в период
с~16 по~17~мая 2026~года. В роли основной языковой модели
(\emph{worker}) для E1--E4 использовалась \code{cerebras:gpt-oss-120b}
с параметром \code{reasoning\_effort=low} для рабочих агентов и
\code{high} для агента-критика. В роли судьи оценки качества
выступала \code{openai:gpt-4.1-mini} с самосогласованностью из трех
независимых вызовов и нулевой температурой выборки.
Шлюз~\code{HumanGateway} (компонент человека в~контуре управления,
от англ. Human-in-the-Loop, HITL) работал в режиме симулятора через
ту же~\code{openai:gpt-4.1-mini} с системной подсказкой, зависящей
от текущей роли. Подтверждающие прогоны (\code{confirmation\_e3},
\code{confirmation\_e4}) выполнены с заменой рабочего агента
на~\code{cerebras:qwen-3-235b-a22b-instruct-2507} (семейство Alibaba
Qwen) при сохранении остального стека неизменным.

Сводные объемные характеристики экспериментальной серии приведены
в~таблице~\ref{tab:exp:totals}. Совокупный объем составил
5175~прогонов и~суммарную стоимость около~\$87~--- существенно ниже
планового бюджета в~\$200~благодаря низкой удельной стоимости
вызовов на~инфраструктуре~Cerebras (от англ. Wafer-Scale Engine,
WSE). Во~всех финальных грид-прогонах доля отказов нулевая.

\begin{table}[!htbp]
\caption{Объемные характеристики экспериментальной серии}
\label{tab:exp:totals}
\small
\begin{tabularx}{\textwidth}{l c c c X}
\toprule
\textbf{Этап} & \textbf{Прогоны} & \textbf{Стоимость, \$} &
\textbf{Время} & \textbf{Назначение} \\
\midrule
E1   & 900  & 12{,}62 & 3{,}7~ч &
   Базовая линия статических топологий, RQ1. \\
E2   & 2295 & 42{,}37 & 13{,}9~ч &
   Влияние роли человека по матрице (топология~$\times$~роль), RQ3. \\
E3   & 540  & 10{,}12 & 12{,}0~ч &
   Адаптивная мета-топология в~трех режимах маршрутизации, RQ2. \\
E4   & 540  & 6{,}38  & 2{,}1~ч &
   Адаптивная роль на~уже-адаптивной топологии, RQ4. \\
\code{conf\_e3} & 360 & 6{,}50  & 1{,}0~ч &
   Кросс-семейная валидация E3 на~семействе Qwen. \\
\code{conf\_e4} & 540 & 9{,}40  & 1{,}3~ч &
   Кросс-семейная валидация E4 на~семействе Qwen. \\
\midrule
\textbf{Итого}  & \textbf{5175}  & \textbf{87{,}39} &
   \textbf{$\approx$34~ч машинного времени} & --- \\
\bottomrule
\end{tabularx}
\end{table}

\subsection{Эксперимент E1 — базовая линия статических топологий}
\label{sec:exp:e1}

Эксперимент~E1 устанавливает верхнюю границу качества для пяти
статических топологий на~четырех корпусах задач~--- программирование
(HumanEval), математические рассуждения (GSM8K), творческая
генерация (CommonGen) и анализ данных (InfiAgent-DABench), далее
DABench,~--- отвечая на~исследовательский вопрос~RQ1. Сетка прогона
составила $5~\text{топологий} \times 4~\text{корпуса} \times
15~\text{задач} \times 3~\text{зерна}~=~900$ прогонов; роль человека
отключена (\code{human.role=none}). Результаты усреднения по~ячейкам
(топология~$\times$~корпус) приведены в~таблице~\ref{tab:exp:e1q}.

\begin{table}[!htbp]
\caption{Среднее качество $q$ статических топологий по корпусам
(жирным выделен победитель в столбце), E1}
\label{tab:exp:e1q}
\small
\begin{tabularx}{\textwidth}{l c c c c c}
\toprule
\textbf{Топология} & \textbf{HumanEval} & \textbf{GSM8K} &
\textbf{CommonGen} & \textbf{DABench} & \textbf{Среднее} \\
\midrule
Chain        & \textbf{0{,}933} & 0{,}911 & 0{,}526 & 0{,}333 & 0{,}676 \\
Star         & 0{,}911 & \textbf{1{,}000} & 0{,}427 & 0{,}282 & 0{,}655 \\
Debate       & 0{,}644 & 0{,}933 & \textbf{0{,}584} & \textbf{0{,}348} & 0{,}627 \\
Hierarchical & 0{,}889 & 0{,}978 & 0{,}486 & 0{,}267 & 0{,}655 \\
Mesh         & 0{,}667 & 0{,}867 & 0{,}397 & 0{,}319 & 0{,}562 \\
\midrule
\textit{Среднее по победителям} & \textit{0{,}933} & \textit{1{,}000}
   & \textit{0{,}584} & \textit{0{,}348} & \textbf{\textit{0{,}716}} \\
\bottomrule
\end{tabularx}
\end{table}

\textbf{Содержательный итог.} Победители на четырех корпусах
различаются~--- chain на HumanEval, star на~GSM8K и debate
на~CommonGen и~DABench. Единой топологии, доминирующей на всех
типах задач, нет. Лучшая глобальная статика~--- chain
со~средним~$q=0{,}676$~--- проигрывает идеальному выбору
\emph{под задачу} (среднее по победителям, $q=0{,}716$) около
четырех процентных пунктов. Этот разрыв задает~\emph{class-level
upper bound} для последующего эксперимента~E3.

Обнаружено два дополнительных явления. Во-первых,
топология~Mesh~деградирует на программировании и творческой
генерации до~$q\le0{,}40$ из-за вырождения голосования
по~совпадению строк. Во-вторых, на корпусе~DABench все статические
топологии демонстрируют~$q\le0{,}35$ при доле полностью неудачных
прогонов~62--71\,\%~--- проявление, которое в~разделе~\ref{sec:exp:conf}
будет переинтерпретировано как специфика рабочего агента,
а не корпуса.

\textbf{Ответ на~RQ1.} Универсального статического победителя
не существует. Оптимальный выбор топологии~--- функция класса
задач, что мотивирует ввод адаптивной мета-топологии в~E3.

\subsection{Эксперимент E2 — роль человека (RQ3)}
\label{sec:exp:e2}

Эксперимент~E2 измеряет влияние роли человека в графе агентов
поверх статических топологий. Сетка построена как декартово
произведение $5~\text{топологий} \times 5~\text{ролей} \times
4~\text{корпуса} \times 15~\text{задач} \times 3~\text{зерна}$
с~двумя фильтрами~--- по~per-task top-3 топологий из~E1 и~по~трем
структурно бессмысленным парам (Debate~$\times$~Coordinator,
Chain~$\times$~Peer, Hierarchical~$\times$~Peer). Итоговый объем
составил~2295~прогонов. Полная матрица победителей роли по~корпусам
приведена в~таблице~\ref{tab:exp:e2roles}.

\begin{table}[!htbp]
\caption{Победитель среди пяти ролей человека по каждому корпусу
с~маргинализацией по~топологиям, E2}
\label{tab:exp:e2roles}
\small
\begin{tabularx}{\textwidth}{l X c c c}
\toprule
\textbf{Корпус} & \textbf{Победившая роль} &
\textbf{$q$ победителя} & \textbf{$n$ ячеек} &
\textbf{$\Delta q$ vs~E1 best static} \\
\midrule
HumanEval & Coordinator & 0{,}948 & 135 & $+0{,}015$ \\
GSM8K     & Monitor     & 0{,}963 & 135 & $-0{,}037$ \\
CommonGen & Peer        & 0{,}643 & 45  & $+0{,}059$ \\
DABench   & Peer        & 0{,}563 & 90  & $+0{,}215$ \\
\bottomrule
\end{tabularx}
\end{table}

Совокупный лифт от~подключения человека~(\emph{HITL lift}),
рассчитанный на~пересечении ячеек~E1~и~E2, составил~$+0{,}033$
по~среднему качеству ($+4{,}8\,\%$ относительного прироста).
Распределение лифта по~ячейкам неоднородно. Положительный лифт
концентрируется на~трех конфигурациях~---
Hierarchical~$\times$~CommonGen ($+0{,}107$),
Chain~$\times$~DABench ($+0{,}087$) и
Mesh~$\times$~DABench ($+0{,}052$); отрицательный лифт наблюдается
только на~Star~$\times$~GSM8K ($-0{,}062$), где топология уже
близка к~потолку качества.

\textbf{Ответ на~RQ3.} Подключение человека дает скромное
улучшение среднего качества при существенно различном эффекте
по~ячейкам. Никакая единая роль не выигрывает на~всех корпусах~---
ключевой аргумент против фиксированного выбора роли и в~пользу
адаптивной маршрутизации, проверяемой в~E4.

\subsection{Эксперимент E3 — адаптивная мета-топология (RQ2)}
\label{sec:exp:e3}

Эксперимент~E3 непосредственно отвечает на~центральный
исследовательский вопрос~RQ2. На~единой адаптивной мета-топологии
сравниваются три режима маршрутизатора~---
\code{rule}~(детерминированное правиловое решение),
\code{llm}~(вызов языковой модели с~парсингом по~Pydantic) и
\code{oracle}~(чтение таблицы~\code{best per task} из~файла,
построенного по~результатам~E1). Каждый режим прогнан в~сетке
$4~\text{корпуса} \times 15~\text{задач} \times 3~\text{зерна}=180$
прогонов; фиксирована роль человека~Reviewer. Сводные показатели
по~режимам приведены в~таблице~\ref{tab:exp:e3overall}.

\begin{table}[!htbp]
\caption{Сводные показатели трех режимов маршрутизатора
топологий, E3}
\label{tab:exp:e3overall}
\small
\begin{tabularx}{\textwidth}{l c c c c X}
\toprule
\textbf{Режим} & \textbf{Среднее $q$} & \textbf{Стоимость, \$} &
\textbf{Время, c} & \textbf{Итераций} & \textbf{Комментарий} \\
\midrule
oracle & 0{,}697 & 0{,}0227 & 1077 & 6{,}27 &
   Верхняя граница, знание корпуса. \\
rule   & 0{,}645 & 0{,}0236 & 928  & 6{,}04 &
   Правиловый, без вызовов языковой модели. \\
llm    & 0{,}631 & \textbf{0{,}0099} & \textbf{881}  & 6{,}10 &
   Вызов языковой модели на каждый переход. \\
\midrule
\textit{Best static (E1)} & \textit{0{,}716} & \textit{0{,}014} &
   \textit{448} & \textit{4{,}35} &
   \textit{Идеальный выбор статики под корпус.} \\
\bottomrule
\end{tabularx}
\end{table}

Численная картина допускает двойственную интерпретацию.
По~абсолютному качеству все три режима маршрутизатора уступают
лучшей статической конфигурации под корпус~--- оракул отстает
на~$1{,}9$ процентных пункта, правиловой~--- на~$7{,}2$,
маршрутизатор на основе языковой модели~--- на~$8{,}5$. Причина~---
организационная накладная стоимость адаптивной обвязки
(фазы~$\times$~подграфы~$\times$~маршрутизатор), которая не
окупается на~хорошо-решаемых корпусах (GSM8K, HumanEval).
Единственный корпус с~положительным абсолютным эффектом адаптации~---
CommonGen, где оракул дает~$+0{,}07$ относительно лучшей статики.

Сравнение трех режимов маршрутизатора между собой~---
\emph{внутри} адаптивного класса~--- выявляет дополнительный
неосновной эффект. Маршрутизатор на~основе языковой модели
проигрывает правиловому всего~$1{,}3$ процентных пункта по~качеству,
но~стоит в~$2{,}38$~раза дешевле. Соотношение стоимостей удерживается
на~всех четырех корпусах в~диапазоне $2{,}12{-}3{,}00$, что
указывает на~\emph{структурный} эффект~--- маршрутизатор на основе
языковой модели систематически выбирает более дешевые подтопологии
(например, Hierarchical с~\code{max\_rounds=2}) против более
прожорливых решений правиловой таблицы. Этот сопутствующий эффект
не отвечает на~RQ2, но представляет самостоятельный интерес и
обсуждается отдельно в~разделе~\ref{sec:exp:conf}.

\textbf{Ответ на~RQ2.} Адаптивная мета-топология не дает значимого
прироста качества и не дает экономии стоимости по~сравнению с~лучшей
статической конфигурацией под корпус. В~качестве референса для
сравнения используется не одна фиксированная статика, а~оракул-выбор
наилучшей статической топологии для каждого корпуса (среднее
качество $0{,}716$). Оракул адаптивной обвязки достигает $0{,}697$
и отстает от~этого референса на~$1{,}9$ процентных пункта; правиловой
и языковой маршрутизаторы~--- на~$7{,}2$ и~$8{,}5$~процентных пункта.
Гипотеза~RQ2 эмпирически не подтверждается. Обнаруженная рядом
сопутствующая находка~--- двукратная экономия стоимости
маршрутизатора на~основе языковой модели внутри адаптивного
класса~--- не отвечает на~RQ2 в~его исходной постановке (сравнение
адаптивной и статической конфигураций), и при кросс-семейной
валидации на~семействе Qwen, где правиловой маршрутизатор Парето-
доминирует над маршрутизатором на~основе языковой модели, эта
находка также не воспроизводится (см.~\ref{sec:exp:conf}).

\subsection{Эксперимент E4 — адаптивная роль (RQ4)}
\label{sec:exp:e4}

Эксперимент~E4 размещает второй слой адаптивности~--- маршрутизатор
ролей~--- поверх уже-адаптивной топологии. На~основе результатов~E3
зафиксирован выбор~\code{topology\_router=llm} как
Парето-оптимальной точки адаптивного класса. Сметаемая переменная~---
\code{role\_router}~$\in$~\{\code{fixed}, \code{rule}, \code{llm}\};
сетка $3 \times 4 \times 15 \times 3=540$ прогонов. Сводка
по~режимам~--- в~таблице~\ref{tab:exp:e4overall}.

\begin{table}[!htbp]
\caption{Сводные показатели трех режимов маршрутизатора ролей
поверх адаптивной мета-топологии, E4}
\label{tab:exp:e4overall}
\small
\begin{tabularx}{\textwidth}{l c c c c}
\toprule
\textbf{Режим} & \textbf{Среднее $q$} & \textbf{$\Delta q$ vs
   \code{fixed}} & \textbf{Стоимость, \$} & \textbf{Время, c} \\
\midrule
fixed & 0{,}653 & ---       & 0{,}01186 & 381 \\
rule  & \textbf{0{,}680} & $+0{,}027$ & 0{,}01193 & 356 \\
llm   & 0{,}643 & $-0{,}010$ & 0{,}01168 & 285 \\
\bottomrule
\end{tabularx}
\end{table}

Маршрутизатор на~основе правил формально увеличивает среднее
качество на~$2{,}7$~процентных пункта над фиксированной ролью
\code{reviewer}. Направление совпадает на~четырех корпусах из~четырех
при~оговорке о~малом размере выборки и одном семействе весов
рабочего агента; драйвер прироста~--- HumanEval ($+0{,}067$).
Стоимость трех режимов укладывается в~коридор шириной~$0{,}85\,\%$~---
дополнительные стоимостные издержки адаптивной маршрутизации роли
незначимы.

Качественный анализ распределения ролей показал, что правиловой
маршрутизатор фактически коллапсирует в~единственную роль~---
координатор активируется в~$98{,}3\,\%$ переходов. То~есть
адаптации роли по~фазам в~текущей реализации не~происходит~---
фактически реализуется почти-статическая политика с~единственной
ролью. E4 переоткрывает результат~E2~\enquote{Coordinator выигрывает
на HumanEval} в~новой обертке. Подробная интерпретация
этого механизма дана в~разделе~\ref{sec:analysis:rq}.

\textbf{Ответ на~RQ4.} Гипотеза о~превосходстве совместной адаптации
топологии и роли над статическими комбинациями эмпирически
не подтверждается. На~исходной семье моделей правиловая
маршрутизация роли дает направленный прирост качества $+0{,}027$
при~нулевой дополнительной стоимости, однако
бутстрап-доверительный интервал $[-0{,}055,\, +0{,}110]$ покрывает
нуль, и~строгий вывод о~значимости прироста не делается. Сопутствующая
находка~--- направление эффекта совпадает на~всех четырех корпусах
из~четырех~--- при~кросс-семейной валидации~(\ref{sec:exp:conf})
также не подтверждается~--- на~семействе Qwen знак эффекта
становится отрицательным.

\subsection{Подтверждающие прогоны — кросс-семейная валидация}
\label{sec:exp:conf}

Два подтверждающих прогона~--- \code{confirmation\_e3}
и~\code{confirmation\_e4}~--- повторяют экспериментальную сетку
E3~и~E4 с~единственной измененной переменной~--- семейством
рабочей языковой модели. Вместо~\code{gpt-oss-120b} использована
\code{qwen-3-235b-a22b-instruct-2507} (семейство Alibaba Qwen).
Все остальные компоненты (судья, шлюз~HITL, инфраструктура,
адаптивная обвязка, защитные правила, конфигурация порогов) сохранены
неизменными. Цель~--- проверить переносимость основных выводов
E3~и~E4 между семействами весов.

Результаты~\code{confirmation\_e3} (\code{topology\_router}~$\in$~\{rule, llm\})
приведены в~таблице~\ref{tab:exp:confe3}. Соотношение стоимостей
правилового маршрутизатора к~маршрутизатору на основе языковой модели
схлопывается с~$2{,}38$ на~исходном семействе до~$1{,}11$~на~семействе
Qwen, а~разрыв качества увеличивается с~$+1{,}3$ процентных пункта
до~$+30{,}6$ в~пользу правилового режима.

\begin{table}[!htbp]
\caption{Сравнение двух режимов маршрутизатора топологий между
семействами весов, \code{confirmation\_e3}}
\label{tab:exp:confe3}
\small
\begin{tabularx}{\textwidth}{l c c c c X}
\toprule
\textbf{Семейство} & \textbf{$q_{\text{rule}}$} &
\textbf{$q_{\text{llm}}$} & \textbf{$\Delta q$} &
\textbf{$c_{\text{rule}}/c_{\text{llm}}$} &
\textbf{Парето-победитель} \\
\midrule
gpt-oss-120b (E3)      & 0{,}645 & 0{,}631 & $+0{,}013$ & $2{,}38$ &
   \code{llm} (по стоимости) \\
qwen-3-235b (conf\_e3) & 0{,}866 & 0{,}560 & \textbf{$+0{,}307$} &
   $1{,}11$ & \code{rule} (по обеим осям) \\
\bottomrule
\end{tabularx}
\end{table}

Аналогичная картина для~\code{confirmation\_e4} (см.~таблицу
\ref{tab:exp:confe4}). На исходном семействе моделей правиловой
маршрутизатор роли превосходил фиксированную конфигурацию
на~$+0{,}027$; на~семействе Qwen он, напротив, уступает
ей~на~$-0{,}081$~--- знак эффекта изменился.

\begin{table}[!htbp]
\caption{Сравнение трех режимов маршрутизатора ролей между
семействами весов, \code{confirmation\_e4}}
\label{tab:exp:confe4}
\small
\begin{tabularx}{\textwidth}{l c c c c}
\toprule
\textbf{Режим} & \textbf{$q$ на gpt-oss} & \textbf{$q$ на qwen} &
\textbf{$\Delta$ vs~\code{fixed} gpt-oss} & \textbf{$\Delta$ vs~\code{fixed} qwen} \\
\midrule
fixed & 0{,}653 & 0{,}582 & ---       & --- \\
rule  & 0{,}680 & 0{,}501 & $+0{,}027$ & \textbf{$-0{,}081$} \\
llm   & 0{,}643 & 0{,}528 & $-0{,}010$ & $-0{,}054$ \\
\bottomrule
\end{tabularx}
\end{table}

Качественный анализ выявил дополнительный эффект~--- сдвиг
распределения ролей при~неизменной таблице~\code{DEFAULT\_ROLE\_TABLE}.
На~исходной семье правиловой маршрутизатор выбирал~\emph{координатор}
в~$98\,\%$ переходов; на~семье~Qwen тот же маршрутизатор с~той же
таблицей выбирает~\emph{равноправного участника} в~$94\,\%$
переходов. Иначе говоря, политика, выглядящая независимой от модели,
фактически зависит от нее через входы~--- фазовые классификаторы
получают разные сигналы от моделей разных семейств и относят те же
ситуации к~разным фазам.

\textbf{Дополнительная находка~--- DABench.} На~исходном семействе
корпус~DABench демонстрировал среднее качество~$0{,}15{-}0{,}24$ во~всех
экспериментах E1--E4; на~семье~Qwen те же конфигурации дают
качество $0{,}84{-}0{,}93$. Это согласуется с~гипотезой о~том, что
DABench измеряет компетенции, специфичные для~семейства весов
рабочего агента, а не является~\enquote{сломанным корпусом};
численный контраст между семействами~--- около 70 процентных пунктов
в~абсолютном выражении~--- исключает корпус как источник сквозного
смещения.

\textbf{Ограничение интерпретации.} На~исходном семействе работники
запускались с~параметром~\code{reasoning\_effort=low}; программный
интерфейс~Qwen этот параметр не~принимает. Часть наблюдаемых
кросс-семейных различий смешана с~различием в~глубине рассуждений
работника и не отделяется от~эффекта семейства весов без
дополнительного контроля; это явное ограничение, требующее
повторной серии при появлении соответствующей опции
в~интерфейсе~Qwen.

\subsection{Бутстрап-доверительные интервалы}
\label{sec:exp:bootstrap}

Для оценки устойчивости заявленных эффектов проведен непараметрический
бутстрап с~десятью тысячами итераций (метод процентилей,
\code{seed=42}). Сводные интервалы для ключевых разностей и
отношений~--- в~таблице~\ref{tab:exp:bootstrap}.

\begin{table}[!htbp]
\caption{Бутстрап-доверительные интервалы $95\,\%$ для ключевых
разностей и отношений}
\label{tab:exp:bootstrap}
\small
\begin{tabularx}{\textwidth}{X c c c}
\toprule
\textbf{Сравнение} & \textbf{Точечная оценка} &
\textbf{Интервал~$95\,\%$} & \textbf{Содержит границу?} \\
\midrule
E3 rule${-}$llm, $q$ (gpt-oss)            & $+0{,}013$ &
   $[-0{,}072,\, +0{,}097]$ & да (нуль) \\
E3 oracle${-}$rule, $q$ (gpt-oss)         & $+0{,}053$ &
   $[-0{,}033,\, +0{,}137]$ & да (нуль) \\
E3 $c_{\text{rule}}/c_{\text{llm}}$       & $2{,}38$   &
   $[1{,}94,\, 2{,}94]$     & \textbf{нет (единицу)} \\
E4 rule${-}$fixed, $q$ (gpt-oss)          & $+0{,}027$ &
   $[-0{,}055,\, +0{,}110]$ & да (нуль) \\
E4 rule${-}$fixed (HumanEval, gpt-oss)    & $+0{,}067$ &
   $[-0{,}022,\, +0{,}156]$ & да (нуль) \\
\code{conf\_e3} rule${-}$llm, $q$ (qwen)  & $+0{,}307$ &
   \textbf{$[+0{,}233,\, +0{,}378]$} & \textbf{нет (нуль)} \\
\code{conf\_e3} $c_{\text{rule}}/c_{\text{llm}}$ (qwen) & $1{,}11$ &
   $[0{,}98,\, 1{,}25]$     & да (единицу) \\
\code{conf\_e4} rule${-}$fixed, $q$ (qwen) & $-0{,}081$ &
   $[-0{,}173,\, +0{,}010]$ & пограничное \\
\bottomrule
\end{tabularx}
\end{table}

Интервалы позволяют отделить устойчивые эффекты от направленных
точечных оценок. Устойчивых эффектов в~работе три. Во-первых,
двукратная экономия стоимости маршрутизатора на основе языковой
модели на~исходной семье~--- интервал отношения стоимостей
$[1{,}94,\, 2{,}94]$ никогда не пересекает~единицу. Во-вторых,
изменение знака эффекта качества правилового маршрутизатора при
смене семьи весов~--- интервал разности качеств на~семье~Qwen
$[+0{,}233,\, +0{,}378]$ устойчиво исключает нуль. В-третьих,
схлопывание соотношения стоимостей при смене семьи~--- интервал
$[0{,}98,\, 1{,}25]$ покрывает единицу, что указывает на
структурный характер схлопывания. Остальные точечные оценки
направленные, но не устойчивы~--- бутстрап-интервалы включают
нуль, и эффекты в строгом смысле статистически неотличимы от
нулевого.

\subsection{Абляционное исследование}
\label{sec:exp:ablation}

Поскольку структура воронки от~E1 до~E4 уже устроена как
последовательное~\emph{добавление} компонентов~--- сначала статические
топологии, затем человек, затем адаптивная топология, затем
адаптивная роль~--- сама воронка играет роль абляционного исследования.
Сводка эффектов компонентов приведена в~таблице~\ref{tab:exp:ablation}.

\begin{table}[!htbp]
\caption{Абляционная сводка~--- эффект добавления компонента
относительно базовой линии}
\label{tab:exp:ablation}
\small
\begin{tabularx}{\textwidth}{l c c X}
\toprule
\textbf{Компонент} & \textbf{$\Delta q$} & \textbf{$\Delta c$} &
\textbf{Источник эффекта} \\
\midrule
Best static per task (E1)            & ---       & ---       &
   Базовая линия. \\
+ HITL фиксированной роли (E2)       & $+0{,}033$ & $+0{,}004$ &
   Подключение симулятора человека. \\
+ Адаптивная мета-топология, oracle (E3)
                                     & $-0{,}019$ & $+0{,}008$ &
   Накладные расходы маршрутизации без явного выигрыша. \\
+ Адаптивная роль, rule (E4 vs E3 llm)
                                     & $+0{,}049$ & $\approx 0$ &
   Поверх E3 llm; скрытое смещение к роли~\emph{coordinator}. \\
+ Смена семейства модели (\code{conf\_e4} vs E4)
                                     & $-0{,}108$ & $+0{,}005$ &
   Падение по знаку и магнитуде~--- переносимость не подтверждена. \\
\bottomrule
\end{tabularx}
\end{table}

Содержательно абляция показывает, что только два компонента~---
подключение человека (E2~vs~E1) и фиксация роли~\emph{coordinator}
через правиловую таблицу (E4~vs~E3)~--- дают положительное
приращение качества; адаптивная топология сама по себе ухудшает
качество относительно идеально-настроенной статики; смена семейства
весов меняет знак эффекта роли. Кросс-семейная неустойчивость
ставит под вопрос интерпретацию положительных приращений как
универсальных, а не специфичных для исходной семьи моделей.

\subsection{Выводы по главе}
\label{sec:exp:summary}

Результаты экспериментальной серии сводятся к~четырем содержательным
итогам. Первый~--- класс задач определяет оптимальную статическую
топологию (RQ1~подтверждается). Второй~--- адаптивная мета-топология
не превосходит лучшую статическую конфигурацию под корпус
ни~по~качеству, ни~по~стоимости (RQ2~не~подтверждается); рядом
обнаружен неосновной структурный эффект~--- двукратная экономия
стоимости маршрутизатора на~основе языковой модели внутри
адаптивного класса~--- который при~кросс-семейной валидации также
не подтверждается. Третий~--- подключение человека дает скромный
направленный прирост качества при~существенно разном эффекте
по~ячейкам (RQ3~подтверждается). Четвертый~--- совместная адаптация
топологии и роли не превосходит фиксированные комбинации
(RQ4~не~подтверждается); направление эффекта совпадает
на~четырех корпусах из~четырех в~исходной семье моделей, но при
кросс-семейной валидации знак меняется и эффект исчезает.

Единственный устойчивый кросс-семейный вывод~--- зависимость
оптимальной политики маршрутизации от~семейства весов рабочего
агента. Его содержательная интерпретация и обсуждение угроз
валидности вынесены в~главу~\ref{sec:analysis}.


\newpage
% Глава 5. Анализ результатов и обсуждение.
% Объем 3--4 страницы. Блок до данных Главы 4.
%
% Структура
%   5.1 Ответы на исследовательские вопросы RQ1--RQ4
%   5.2 Главный нарратив исследования
%   5.3 Сопоставление полученных результатов с гипотезами
%   5.4 Угрозы валидности и ограничения

\section{Анализ результатов и обсуждение}
\label{sec:analysis}

\TODO{Этап 8. Блок до завершения Главы 4.}


%%% ===== Авторский вклад (без \newpage --- идет следом за Главой 5) ===== %%%
% Авторский вклад --- отдельный раздел, требование методички, п. 3.5.
% Объем 0.5--1 страница.

\section*{Авторский вклад}
\label{sec:contribution}
\addcontentsline{toc}{section}{Авторский вклад}

\TODO{Этап 8 заключительный. Три пункта вклада.
1. Программная инфраструктура мульти-агентного фреймворка
   --- пять статичных топологий, адаптивная мета-топология,
   шлюз HITL, три режима маршрутизации, leave-one-out оракул.
2. Таксономия метрик для адаптивных мульти-агентных систем.
3. Эмпирические результаты на четырех классах задач.}


%%% ===== Заключение ===== %%%
\newpage
% Заключение. Версия 2 --- после получения данных Главы 4.
% Объем 2 страницы.
% Раздел «Направления дальнейших исследований» --- связным текстом,
% минимум 10 строк (правки друга-ревьюера).

\section*{Заключение}
\addcontentsline{toc}{section}{Заключение}

В работе проведена количественная оценка влияния выбора
коммуникационной топологии, механизма ее адаптации и~позиции
человека на~эффективность решения задач разных классов MAS LLM.
Цель достигнута~--- получены количественные характеристики для
пяти статических топологий, адаптивной мета-топологии в~двух
режимах маршрутизации, пяти ролей человека и~их~адаптивной
комбинации; выполнено 4230~завершенных запусков на~четырех корпусах
задач с~суммарной стоимостью~\$105{,}35.

RQ1 (от англ. Research Question) подтвержден~--- универсальной
статики-доминатора нет, победители на~четырех корпусах различаются
(Chain~--- HumanEval, Star~--- GSM8K, Debate~--- CommonGen и~DABench).
RQ3 подтвержден~--- HITL дает прирост $+0{,}033$ при~неоднородности
по~ячейкам.

RQ2 раскладывается на~центральный вывод~--- зависящее от~сложности
задачи Парето-преимущество адаптивной мета-топологии на~нетривиальных
корпусах с~поведением наименьшего риска (Утверждение~1)~---
и~поддерживающее эмпирическое наблюдение об~инверсии Парето-победителя
между развернутыми конфигурациями работника при~переходе с~gpt-oss-120b
на~Qwen-3-235B (Утверждение~2; см.~\ref{sec:analysis:rq}). RQ4
опровергнут с~интерпретацией <<универсальный отрицательный результат>>
(Утверждение~3)~--- адаптивный маршрутизатор роли повсеместно
не~дает положительного прироста поверх адаптивной мета-топологии;
эффект не~реализационный, а~в~самой гипотезе двухуровневой адаптации.

Ключевой результат работы~--- количественная демонстрация ценности
адаптивного дизайна на~нетривиальных задачах. На~корпусах, где
единичная статическая топология имеет реальные ограничения
(CommonGen, DABench), адаптивная мета-топология обеспечивает
Парето-выигрыш и~поведение наименьшего риска по~трем независимым
осям (наихудший случай качества, дисперсия качества, разброс
стоимости), сохраняя при~этом структурные свойства дизайна~--- единый
конфиг развертывания, интеграцию HITL-сигнала и~фазовое переключение
под-топологий. Результат получен при~контролируемых условиях
одного семейства работника и~устойчив к~основным сценариям
конфаунда. Поддерживающее наблюдение по~инверсии Парето-победителя
мотивирует практику кросс-конфигурационной валидации настроек MAS
перед~переносом на~новую развернутую конфигурацию работника;
механизм инверсии (семейство весов, глубина рассуждений, поддержка
структурированного вывода) в~работе не~устанавливается. Универсальный
отрицательный результат по~RQ4 дает конкретное проектное
ограничение~--- адаптивность на~двух уровнях одновременно (топология
и~роль) избыточна, достаточно одного уровня (топологии).

Ограничения~--- $n=45$ запусков на~ячейку $(\text{корпус},\text{режим})$
в~E3 и~E4 (на~этом объеме устойчиво подтверждены три эффекта);
кросс-конфигурационная валидация на~двух конфигурациях рабочего агента;
наблюдаемая инверсия Парето-победителя смешана сразу с~несколькими
переменными (семейство весов, глубина рассуждений, поддержка
структурированного вывода) и~в~работе не~разделяется; DABench
проявляет специфичное для~модели ограничение на~gpt-oss-120b
(см.~\ref{sec:exp:conf}); человек моделируется языковой моделью;
корпус не~покрывает ряд классов (диалог, перевод, обобщение длинных
текстов); судья~--- единственная модель из~единственного семейства.

\subsection*{Направления дальнейших исследований}

Прямое продолжение работы образует четыре взаимодополняющих направления.
Первое~--- расширение кросс-семейной валидации на~три и~более семейства
весов рабочего агента (Anthropic Claude, Google Gemini, Meta LLaMA-4)
с~приведением параметра глубины рассуждений к~единой шкале, что устранит
смешение эффектов семейства и~режима работы агента и~позволит проверить,
является ли наблюдаемая кросс-семейная неустойчивость политик
маршрутизации универсальным свойством адаптивных систем или специфичной
для~конкретной пары семейств. Второе~--- валидация симулятора человека
на~реальных участниках по~схеме планируемого эксперимента~E5
с~использованием шкалы~NASA-TLX (от англ. NASA Task Load Index~---
индекс когнитивной нагрузки NASA), латино-квадратного дизайна
на~восьми-двенадцати участниках и~отобранного подмножества
из~двенадцати задач, что позволит численно оценить корреляцию между
симулированной и~фактической оценкой когнитивной нагрузки.
Третье~--- разработка семейно-осознанного маршрутизатора, в~котором
политика выбора топологии и~роли явно учитывает семейство весов
рабочего агента, например, через предварительное дообучение
маршрутизатора на~трассах целевой модели или через переключение между
семейно-специфичными политиками. Четвертое~--- расширение корпуса
задач на~непокрытые классы (диалог, перевод, обобщение длинных
документов, мультимодальные задачи) и~проверка обобщения результатов
воронки на~новые классы.


%%% ===== Библиография ===== %%%
\newpage
% =====================================================================
% Библиографический список.
% Оформление по ГОСТ Р 7.0.5-2008.
% Сортировка строго по алфавиту фамилии первого автора (для записей
% с заглавием в начале --- по первой букве заглавия).
% Минимум 30 источников; в списке 44 наименования.
% =====================================================================

% Защита от переноса внутри слов в заголовке списка литературы:
% русский babel может разорвать "список" → "спи-сок". Длинные слова
% в \refname оборачиваем \mbox-ом.
\renewcommand{\refname}{\mbox{Библиографический} \mbox{список}}
\addcontentsline{toc}{section}{Библиографический список}

\begingroup
\footnotesize
\setlength{\itemsep}{1pt}
\begin{multicols}{2}
\begin{thebibliography}{99}

\bibitem{wang2024}
A survey on large language model based autonomous agents
/ L. Wang, C. Ma, X. Feng [и др.] // Frontiers of Computer Science.
-- 2024. -- Vol. 18, no. 6. -- P. 186345.

\bibitem{agentbench2023}
AgentBench: Evaluating LLMs as agents [Электронный ресурс]
/ X. Liu, H. Yu, H. Zhang [и др.] // arXiv preprint 2308.03688.
-- 2023. -- URL: \url{https://arxiv.org/abs/2308.03688} (дата
обращения 17.05.2026).

\bibitem{agentverse2023}
AgentVerse: Facilitating multi-agent collaboration and exploring
emergent behaviors [Электронный ресурс] / W. Chen, Y. Su, J. Zuo
[и др.] // arXiv preprint 2308.10848. -- 2023.
-- URL: \url{https://arxiv.org/abs/2308.10848} (дата обращения: 17.05.2026).

\bibitem{wu2023}
AutoGen: Enabling next-gen LLM applications via multi-agent
conversation [Электронный ресурс] / Q. Wu, G. Bansal, J. Zhang
[и др.] // arXiv preprint 2308.08155. -- 2023.
-- URL: \url{https://arxiv.org/abs/2308.08155} (дата обращения: 17.05.2026).

\bibitem{benjamini1995}
Benjamini Y. Controlling the false discovery rate: a practical
and powerful approach to multiple testing / Y. Benjamini,
Y. Hochberg // Journal of the Royal Statistical Society.
Series B. -- 1995. -- Vol. 57, no. 1. -- P. 289--300.

\bibitem{camel2023}
CAMEL: Communicative agents for «mind» exploration of large
language model society / G. Li, H. Hammoud, H. Itani, D. Khizbullin,
B. Ghanem // Advances in Neural Information Processing Systems.
-- Red Hook, NY: Curran Associates, 2023. -- Vol. 36.
-- P. 51991--52008.

\bibitem{langchain2023}
Chase H. LangChain: Building applications with LLMs through
composability [Электронный ресурс] / H. Chase. -- 2023.
-- URL: \url{https://github.com/langchain-ai/langchain} (дата
обращения: 17.05.2026).

\bibitem{qian2024}
ChatDev: Communicative agents for software development
/ C. Qian, W. Liu, H. Liu [и др.] // Proceedings of the 62nd
Annual Meeting of the Association for Computational Linguistics
(ACL 2024). -- Stroudsburg, PA: ACL, 2024. -- P. 15174--15186.

\bibitem{frugalgpt2023}
Chen L. FrugalGPT: How to use large language models while
reducing cost and improving performance [Электронный ресурс]
/ L. Chen, M. Zaharia, J. Zou // Transactions on Machine
Learning Research. -- 2024.
-- URL: \url{https://openreview.net/forum?id=cSimKw5p6R} (дата
обращения: 17.05.2026).

\bibitem{christiano2017}
Christiano P. F. Deep reinforcement learning from human
preferences / P. F. Christiano, J. Leike, T. Brown [и др.]
// Advances in Neural Information Processing Systems. -- Red
Hook, NY: Curran Associates, 2017. -- Vol. 30. -- P. 4299--4307.

\bibitem{commongen2020}
CommonGen: A constrained text generation challenge for generative
commonsense reasoning / B. Y. Lin, W. Zhou, M. Shen [и др.]
// Findings of the Association for Computational Linguistics
(EMNLP 2020). -- Stroudsburg, PA: ACL, 2020. -- P. 1823--1840.

\bibitem{cohen1988}
Cohen J. Statistical power analysis for the behavioral sciences
/ J. Cohen. -- 2nd ed. -- Hillsdale, NJ: Lawrence Erlbaum
Associates, 1988. -- 567 p.

\bibitem{pydantic2024}
Colvin S. Pydantic: Data validation using Python type hints
[Электронный ресурс] / S. Colvin [и др.]. -- 2024.
-- URL: \url{https://docs.pydantic.dev/} (дата обращения: 17.05.2026).

\bibitem{efron1979}
Efron B. Bootstrap methods: another look at the jackknife
/ B. Efron // The Annals of Statistics. -- 1979. -- Vol. 7,
no. 1. -- P. 1--26.

\bibitem{humaneval2021}
Evaluating large language models trained on code [Электронный
ресурс] / M. Chen, J. Tworek, H. Jun [и др.]
// arXiv preprint 2107.03374. -- 2021.
-- URL: \url{https://arxiv.org/abs/2107.03374} (дата обращения: 17.05.2026).

\bibitem{zhang2024}
G-Designer: Architecting multi-agent communication topologies
via graph neural networks [Электронный ресурс] / Y. Zhang,
K. Xu, Y. Li, Z. Liu, D. Shen // arXiv preprint 2410.11782.
-- 2024. -- URL: \url{https://arxiv.org/abs/2410.11782} (дата
обращения 17.05.2026).

\bibitem{zhuge2024}
GPTSwarm: Language agents as optimizable graphs / M. Zhuge,
W. Wang, L. Kirsch, F. Faccio, D. Khizbullin, J. Schmidhuber
// Proceedings of the International Conference on Machine
Learning (ICML 2024). -- Cambridge, MA: PMLR, 2024.
-- P. 62556--62583.

\bibitem{nasatlx1988}
Hart S. G. Development of NASA-TLX (Task Load Index): results
of empirical and theoretical research / S. G. Hart,
L. E. Staveland // Human Mental Workload / ed. by P. A. Hancock,
N. Meshkati. -- Amsterdam: North-Holland, 1988. -- P. 139--183.

\bibitem{hellaswag2019}
HellaSwag: Can a machine really finish your sentence?
/ R. Zellers, A. Holtzman, Y. Bisk, A. Farhadi, Y. Choi
// Proceedings of the 57th Annual Meeting of the Association
for Computational Linguistics (ACL 2019). -- Stroudsburg, PA:
ACL, 2019. -- P. 4791--4800.

\bibitem{horvitz1999}
Horvitz E. Principles of Mixed-Initiative User Interfaces
/ E. Horvitz // Proceedings of the SIGCHI Conference on Human
Factors in Computing Systems (CHI '99). -- New York: ACM, 1999.
-- P. 159--166.

\bibitem{du2023}
Improving factuality and reasoning in language models through
multiagent debate [Электронный ресурс] / Y. Du, S. Li,
A. Torralba, J. Tenenbaum, I. Mordatch // arXiv preprint
2305.14325. -- 2023.
-- URL: \url{https://arxiv.org/abs/2305.14325} (дата обращения: 17.05.2026).

\bibitem{dabench2024}
InfiAgent-DABench: Evaluating agents on data analysis tasks
[Электронный ресурс] / X. Hu, Z. Zhao, S. Wei [и др.]
// arXiv preprint 2401.05507. -- 2024.
-- URL: \url{https://arxiv.org/abs/2401.05507} (дата обращения: 17.05.2026).

\bibitem{humanevalplus2023}
Is your code generated by ChatGPT really correct? Rigorous
evaluation of large language models for code generation
/ J. Liu, C. S. Xia, Y. Wang, L. Zhang // Advances in Neural
Information Processing Systems. -- Red Hook, NY: Curran
Associates, 2023. -- Vol. 36. -- P. 21558--21572.

\bibitem{langgraph2024}
LangGraph: A framework for building stateful, multi-actor
applications with LLMs [Электронный ресурс] / LangChain AI.
-- Электрон. дан. -- 2024.
-- URL: \url{https://langchain-ai.github.io/langgraph/} (дата
обращения 17.05.2026).

\bibitem{langsmith2024}
LangSmith: Observability and evaluation for LLM applications
[Электронный ресурс] / LangChain AI. -- Электрон. дан. -- 2024.
-- URL: \url{https://docs.smith.langchain.com/} (дата обращения: 17.05.2026).

\bibitem{rouge2004}
Lin C.-Y. ROUGE: A package for automatic evaluation of summaries
/ C.-Y. Lin // Text Summarization Branches Out: Proceedings
of the ACL-04 Workshop. -- Stroudsburg, PA: ACL, 2004.
-- P. 74--81.

\bibitem{livecodebench2024}
LiveCodeBench: Holistic and contamination-free evaluation of
large language models for code [Электронный ресурс] / N. Jain,
K. Han, A. Gu [и др.] // arXiv preprint 2403.07974. -- 2024.
-- URL: \url{https://arxiv.org/abs/2403.07974} (дата обращения: 17.05.2026).

\bibitem{liu2023}
Liu Z. Dynamic LLM-agent network: an LLM-agent collaboration
framework with agent team optimization [Электронный ресурс]
/ Z. Liu, Y. Zhang, P. Li, Y. Liu, D. Yang // arXiv preprint
2310.02170. -- 2023.
-- URL: \url{https://arxiv.org/abs/2310.02170} (дата обращения: 17.05.2026).

\bibitem{magenticone2024}
Magentic-One: A generalist multi-agent system for solving complex
tasks [Электронный ресурс] / A. Fourney, G. Bansal, H. Mozannar
[и др.] // arXiv preprint 2411.04468. -- 2024.
-- URL: \url{https://arxiv.org/abs/2411.04468} (дата обращения: 17.05.2026).

\bibitem{mathbench2021}
Measuring mathematical problem solving with the MATH dataset
/ D. Hendrycks, C. Burns, S. Kadavath [и др.] // Advances in
Neural Information Processing Systems. -- Red Hook, NY: Curran
Associates, 2021. -- Vol. 34. -- P. 2280--2292.

\bibitem{hong2024}
MetaGPT: Meta programming for a multi-agent collaborative
framework / S. Hong, M. Zhuge, J. Chen [и др.]
// Proceedings of the International Conference on Learning
Representations (ICLR 2024). -- 2024.

\bibitem{moa2024}
Mixture-of-Agents enhances large language model capabilities
/ J. Wang, J. Wang, B. Athiwaratkun, C. Zhang, J. Zou
// Proceedings of the Conference on Language Modeling
(COLM 2024). -- 2024.

\bibitem{mmlupro2024}
MMLU-Pro: A more robust and challenging multi-task language
understanding benchmark [Электронный ресурс] / Y. Wang, X. Ma,
G. Zhang [и др.] // arXiv preprint 2406.01574. -- 2024.
-- URL: \url{https://arxiv.org/abs/2406.01574} (дата обращения: 17.05.2026).

\bibitem{mosqueira2023}
Mosqueira-Rey E. Human-in-the-loop machine learning: a state
of the art / E. Mosqueira-Rey, E. Hern\'{a}ndez-Pereira,
D. Alonso-R\'{i}os, J. Bobes-Bascar\'{a}n,
\'{A}. Fern\'{a}ndez-Leal // Artificial Intelligence Review.
-- 2023. -- Vol. 56, no. 4. -- P. 3005--3054.

\bibitem{yu2024}
NetSafe: Exploring the topological safety of multi-agent
networks against adversarial attacks [Электронный ресурс]
/ M. Yu, S. Wang, G. Zhang, H. Li // arXiv preprint 2410.15686.
-- 2024.
-- URL: \url{https://arxiv.org/abs/2410.15686} (дата обращения: 17.05.2026).

\bibitem{ouyang2022}
Ouyang L. Training language models to follow instructions with
human feedback / L. Ouyang, J. Wu, X. Jiang [и др.]
// Advances in Neural Information Processing Systems. -- Red
Hook, NY: Curran Associates, 2022. -- Vol. 35.
-- P. 27730--27744.

\bibitem{react2022}
ReAct: Synergizing reasoning and acting in language models
/ S. Yao, J. Zhao, D. Yu, N. Du, I. Shafran, K. Narasimhan,
Y. Cao // Proceedings of the International Conference on
Learning Representations (ICLR 2023). -- 2023.

\bibitem{reflexion2023}
Reflexion: Language agents with verbal reinforcement learning
/ N. Shinn, F. Cassano, A. Gopinath, K. Narasimhan, S. Yao
// Advances in Neural Information Processing Systems. -- Red
Hook, NY: Curran Associates, 2023. -- Vol. 36.
-- P. 8634--8652.

\bibitem{routellm2024}
RouteLLM: Learning to route LLMs with preference data
[Электронный ресурс] / I. Ong, A. Almahairi, V. Wu [и др.]
// arXiv preprint 2406.18665. -- 2024.
-- URL: \url{https://arxiv.org/abs/2406.18665} (дата обращения: 17.05.2026).

\bibitem{qian2024scaling}
Scaling large-language-model-based multi-agent collaboration
[Электронный ресурс] / C. Qian, Z. Ye, W. Liu [и др.]
// arXiv preprint 2406.07155. -- 2024.
-- URL: \url{https://arxiv.org/abs/2406.07155} (дата обращения: 17.05.2026).

\bibitem{selfrefine2024}
Self-Refine: Iterative refinement with self-feedback / A. Madaan,
N. Tandon, P. Gupta [и др.] // Advances in Neural Information
Processing Systems. -- Red Hook, NY: Curran Associates, 2023.
-- Vol. 36. -- P. 46534--46594.

\bibitem{liang2024}
Tipping the dominos: Topology-aware multi-hop attacks on
LLM-based multi-agent systems [Электронный ресурс] / R. Liang,
J. Ye, Z. Chen, Q. Zhu // arXiv preprint 2412.04129. -- 2024.
-- URL: \url{https://arxiv.org/abs/2412.04129} (дата обращения: 17.05.2026).

\bibitem{gsm8k2021}
Training verifiers to solve math word problems [Электронный
ресурс] / K. Cobbe, V. Kosaraju, M. Bavarian [и др.]
// arXiv preprint 2110.14168. -- 2021.
-- URL: \url{https://arxiv.org/abs/2110.14168} (дата обращения: 17.05.2026).

\bibitem{tot2023}
Tree of Thoughts: Deliberate problem solving with large language
models / S. Yao, D. Yu, J. Zhao, I. Shafran, T. Griffiths, Y. Cao,
K. Narasimhan // Advances in Neural Information Processing
Systems. -- Red Hook, NY: Curran Associates, 2023. -- Vol. 36.
-- P. 11809--11822.

\end{thebibliography}
\end{multicols}
\endgroup


%%% ===== Приложения ===== %%%
\newpage
\appendix
\renewcommand{\thesection}{\Asbuk{section}}
\titleformat{\section}{\normalfont\Large\bfseries}{Приложение~\thesection}{1em}{}

% =====================================================================
% Приложение А. Функциональная схема программной системы.
% Реализовано в TikZ. Версия 4 — горизонтальный snake-flow на
% альбомной странице, слои представлены крупными блоками с перечнем
% содержимого внутри, без \resizebox.
% =====================================================================

\clearpage
\section{Функциональная схема системы}
\label{app:uml}

\renewcommand{\thefigure}{\thesection.\arabic{figure}}
\setcounter{figure}{0}

В настоящем приложении приведена функциональная схема программной
системы, описанной в главе~\ref{sec:architecture}. Схема изображает
движение управления и данных между основными архитектурными блоками
от поступления входных данных эксперимента до получения результирующих
метрик и журналов. Поток выполнения организован по~двум рядам слева
направо, переход на~второй ряд показан изгибом стрелки справа.
На~рисунке~\ref{fig:appA:scheme} двойная красная рамка отмечает
авторские компоненты~--- слой адаптации, оракул-маршрутизатор и~шлюз
взаимодействия с~человеком.

\begin{sidewaysfigure}[p]
\centering
\begin{tikzpicture}[
  font=\small,
  every node/.style={align=center},
  io/.style={rectangle, draw=black, thick, fill=gray!12,
             text width=3.4cm, minimum height=2.0cm, rounded corners=3pt},
  ctrl/.style={rectangle, draw=black, thick, fill=blue!10,
               text width=3.4cm, minimum height=2.0cm, rounded corners=3pt},
  authctrl/.style={rectangle, draw=red!75!black, very thick, double,
                   double distance=1.5pt, fill=blue!10,
                   text width=3.4cm, minimum height=2.0cm, rounded corners=3pt},
  topo/.style={rectangle, draw=black, thick, fill=orange!20,
               text width=3.4cm, minimum height=2.0cm, rounded corners=3pt},
  agent/.style={rectangle, draw=black, thick, fill=green!18,
                text width=3.4cm, minimum height=2.0cm, rounded corners=3pt},
  infra/.style={rectangle, draw=black, thick, fill=red!12,
                text width=3.4cm, minimum height=2.0cm, rounded corners=3pt},
  authinfra/.style={rectangle, draw=red!75!black, very thick, double,
                    double distance=1.5pt, fill=red!12,
                    text width=3.4cm, minimum height=2.0cm, rounded corners=3pt},
  eval/.style={rectangle, draw=black, thick, fill=violet!15,
               text width=3.4cm, minimum height=2.0cm, rounded corners=3pt},
  autheval/.style={rectangle, draw=red!75!black, very thick, double,
                   double distance=1.5pt, fill=violet!15,
                   text width=3.4cm, minimum height=2.0cm, rounded corners=3pt},
  store/.style={rectangle, draw=black, thick, fill=yellow!22,
                text width=3.4cm, minimum height=2.0cm, rounded corners=3pt},
  flow/.style={-{Latex[length=3mm]}, very thick},
]

% ===== Ряд 1 (сверху) =====
\node[io]            (input)    at (0,   0)  {\textbf{Вход}\\Задача $\mathcal{T}$\\и YAML-конфиг};
\node[ctrl]          (orch)     at (4.0, 0)  {\textbf{Оркестратор}\\\code{atm.experiment}\\инициализация};
\node[authctrl]      (adapt)    at (8.0, 0)  {\textbf{Слой адаптации $\bigstar$}\\PhaseRouter\\TopologyRouter\\SwitchGuards};
\node[topo]          (topo)     at (12.0, 0) {\textbf{Топологии}\\Star, Chain, Mesh,\\Debate, Hierarchical\\(активна одна)};
\node[agent]         (agents)   at (16.0, 0) {\textbf{Агенты}\\Planner, Researcher,\\Executor, Critic,\\Debater, Coordinator};
\node[infra]         (services) at (20.0, 0) {\textbf{Внешние службы}\\Инструменты+Docker\\\code{LLMWrapper}};

% ===== Ряд 2 (снизу, движение справа налево) =====
\node[authinfra]     (human)    at (20.0, -5.0) {\textbf{Шлюз человека $\bigstar$}\\\code{HumanGateway}\\(пять ролей)};
\node[ctrl]          (obs)      at (16.0, -5.0) {\textbf{Наблюдаемость}\\\code{atm.observability}\\захват событий};
\node[eval]          (judge)    at (12.0, -5.0) {\textbf{Слой оценки}\\\code{atm.evaluation}\\судья + NASA-TLX};
\node[autheval]      (oracle)   at (8.0, -5.0)  {\textbf{Оракул $\bigstar$}\\per-task best-of\\$\Delta_{\text{oracle}}$};
\node[store]         (store)    at (4.0, -5.0)  {\textbf{Хранение}\\PostgreSQL 16\\Apache Parquet};
\node[io]            (output)   at (0,   -5.0)  {\textbf{Выход}\\Журналы +\\метрики $q, c, \tau,$\\$N_{\text{switch}}, r_{\text{guard}}$};

% ===== Связи ряда 1 =====
\draw[flow] (input)    -- (orch);
\draw[flow] (orch)     -- (adapt);
\draw[flow] (adapt)    -- (topo);
\draw[flow] (topo)     -- (agents);
\draw[flow] (agents)   -- (services);

% ===== Изгиб с ряда 1 на ряд 2 =====
\draw[flow] (services.east) -- ++(0.6, 0)
                            |- (human.east);

% ===== Связи ряда 2 (справа налево) =====
\draw[flow] (human)    -- (obs);
\draw[flow] (obs)      -- (judge);
\draw[flow] (judge)    -- (oracle);
\draw[flow] (oracle)   -- (store);
\draw[flow] (store)    -- (output);

% ===== Обратная связь: оракул → слой адаптации =====
\draw[flow, dashed, red!70!black]
    (oracle.north) -- ++(0, 1.2)
    node[midway, right=2pt, font=\footnotesize, black]
       {оракульная таблица}
    -| (adapt.south);

% ===== Обратная связь: шлюз человека → слой адаптации =====
\draw[flow, dashed, red!70!black]
    (human.north) -- ++(0, 1.7)
    node[midway, right=2pt, font=\footnotesize, black]
       {HITL-сигналы}
    -| (adapt.south east);

\end{tikzpicture}

\caption{Функциональная схема программной системы. Поток управления
организован по~двум рядам слева направо с~переходом на~второй ряд
справа. Двойная красная рамка обозначает авторские компоненты~---
слой адаптации с~маркером $\bigstar$ (маршрутизаторы фаз
и~топологий с~защитными правилами \code{SwitchGuards}),
шлюз взаимодействия с~человеком и~оракул-маршрутизатор. Штриховые
красные стрелки --- обратные связи: оракульная таблица возвращается
в~маршрутизатор топологий, HITL-сигналы поступают от~человека
к~слою адаптации. Полный перечень заимствованных и~авторских
компонентов приведен в~таблице~\ref{tab:arch:contribution}.}
\label{fig:appA:scheme}
\end{sidewaysfigure}

\subsection*{Пояснение к схеме}

Поток выполнения начинается с поступления конфигурации эксперимента
в формате YAML и описания задачи $\mathcal{T}$ из выбранного
корпуса. Оркестратор эксперимента (\code{atm.experiment})
инициализирует общее состояние графа и передает управление слою
адаптации. Маршрутизатор фаз принимает решение о текущей фазе
решения, передает результат маршрутизатору топологий, который
выбирает активную коммуникационную топологию. В работе одновременно
действует ровно одна из пяти базовых топологий, а адаптивный режим
реализуется через последовательное переключение между ними по
решению маршрутизатора.

Активная топология определяет порядок активации агентов и протокол
обмена сообщениями между ними. Любая из шести ролей агентов
(планировщик, исследователь, исполнитель, критик, дебатер,
координатор) может быть подключена к любой топологии в соответствии
с правилами раздела~\ref{sec:meth:topologies}. Агенты обращаются
за внешними службами --- инструментами и Docker-песочницей для
исполнения кода, оболочкой над языковыми моделями для генерации
текста, шлюзом взаимодействия с человеком.

Все события системы (вызовы языковых моделей, выставление сигналов,
переходы между фазами, переключения топологий, обращения к человеку,
срабатывания защитных правил маршрутизатора) перехватываются
подсистемой наблюдаемости и направляются в два хранилища ---
реляционную базу для метаданных и колоночные файлы Apache Parquet
для объемных журналов. На финальной стадии результаты прогона
становятся доступными для постобработки слою анализа
\code{atm.analysis}.

Обратный канал от шлюза человека к маршрутизатору топологий
обозначает поступление управляющих сигналов от участника-человека
к адаптивному слою. Через этот канал человек, действующий в роли
координатора или наблюдателя, может форсировать переключение
топологии или досрочно завершить фазу проверки. Штриховая связь
от~оракула к~маршрутизатору топологий обеспечивает возврат
оракульной таблицы (per-task best-of) для расчета
$\Delta_{\text{oracle}}$.

\newpage
% =====================================================================
% Приложение В. Глоссарий терминов.
% ~30 терминов. Сгруппированы по тематическим блокам.
% =====================================================================

\section{Глоссарий}
\label{app:glossary}

В приложении приведены определения ключевых терминов работы.
Термины сгруппированы по тематическим блокам --- мульти-агентные
системы и языковые модели, коммуникационные топологии, фазы
и маршрутизаторы, человек в контуре управления, метрики и
инфраструктура, эксперимент и статистика. В пределах каждого
блока термины упорядочены логически от общего к частному.

\subsection*{Мульти-агентные системы и языковые модели}

\begin{description}[leftmargin=*, style=nextline]
\item[\textbf{Большая языковая модель} (от англ. Large Language
   Model, LLM)] Глубокая нейронная сеть с большим числом параметров,
   обученная на текстовом корпусе для предсказания следующего
   токена. В настоящей работе языковые модели используются
   в качестве рабочих агентов, симулятора человека, маршрутизаторов
   и судьи.
\item[\textbf{Мульти-агентная система} (от англ. Multi-Agent
   System, MAS)] Программная система, в которой несколько
   автономных агентов взаимодействуют между собой и с внешней
   средой для совместного решения задач.
\item[\textbf{Агент}] Программная сущность в мульти-агентной
   системе с фиксированной ролью, системной подсказкой, набором
   доступных инструментов и собственным внутренним состоянием.
\item[\textbf{Системная подсказка}] Постоянная инструкция, задающая
   агенту его роль, ограничения и стиль ответов. Не изменяется
   в пределах одного запуска.
\item[\textbf{Скрэтчпад} (от англ. scratchpad)] Внутренний
   локальный контекст агента, в котором хранится история
   сообщений и промежуточных рассуждений. Является приватным
   для агента в отличие от общего состояния графа. Организован
   по политике скользящего окна с автоматическим сжатием
   при переполнении.
\item[\textbf{Рабочий агент} (от англ. worker)] Агент, непосредственно
   участвующий в решении задачи, --- планировщик, исследователь,
   исполнитель, критик или дебатер. Противопоставляется
   управляющим компонентам --- координатору, маршрутизатору,
   судье оценки.
\item[\textbf{Судья оценки качества} (от англ. LLM-as-judge)]
   Внешняя языковая модель из другого семейства весов, вызываемая
   для оценки качества финального ответа системы. Использует
   несколько независимых вызовов с агрегацией по правилу
   самосогласованности.
\item[\textbf{Симулятор человека}] Языковая модель в роли участника
   мульти-агентной системы, имитирующая поведение человека в
   контуре управления. Управляется системной подсказкой,
   зависящей от заданной роли --- координатор, рецензент, судья,
   равноправный участник или наблюдатель.
\item[\textbf{Состояние графа} (от англ. graph state)] Общий
   контейнер, объединяющий контексты всех агентов, текущую
   фазу решения, активную топологию, журнал переходов и
   накопленные сигналы. Передается между узлами при исполнении
   и подвергается редукции при параллельном исполнении.
\item[\textbf{Бюджет прогона} (от англ. budget)] Верхний предел
   стоимости вызовов языковых моделей, выделенный на один
   прогон. Реализован в виде трехуровневого ограничителя ---
   на отдельный вызов, на полный прогон и на грид-эксперимент
   в целом.
\end{description}

\subsection*{Коммуникационные топологии}

\begin{description}[leftmargin=*, style=nextline]
\item[\textbf{Коммуникационная топология}] Направленный граф,
   описывающий допустимые маршруты обмена сообщениями между
   агентами мульти-агентной системы и порядок их активации.
\item[\textbf{Топология Star} (звезда)] Конфигурация, в которой
   центральный координатор последовательно вызывает специализированных
   работников и агрегирует их ответы.
\item[\textbf{Топология Chain} (цепочка)] Линейный конвейер из
   трех агентов с обратными связями для доработки.
\item[\textbf{Топология Mesh} (полносвязная)] Конфигурация с общей
   шиной сообщений, в которой все агенты обмениваются сообщениями
   со всеми по принципу циклической очереди.
\item[\textbf{Топология Debate} (дебатная)] Конфигурация, в которой
   два оппонента параллельно генерируют альтернативные решения,
   а судья выбирает победителя или синтезирует ответ.
\item[\textbf{Топология Hierarchical} (иерархическая)] Двухуровневая
   конфигурация из координатора верхнего уровня и двух подкоманд,
   работающих параллельно.
\item[\textbf{Адаптивная мета-топология}] Конфигурация, в которой
   активная базовая топология выбирается на каждой итерации
   маршрутизатором в зависимости от текущей фазы решения и
   наблюдаемых сигналов агентов. Авторская разработка настоящей
   работы.
\end{description}

\subsection*{Фазы и маршрутизаторы}

\begin{description}[leftmargin=*, style=nextline]
\item[\textbf{Фаза решения задачи}] Состояние конечного автомата
   процесса решения. Множество фаз содержит четыре значения ---
   планирование, исполнение, проверка и завершение.
\item[\textbf{Конечный автомат} (от англ. Finite State Machine,
   FSM)] Математическая модель вычислений с конечным множеством
   состояний и переходов между ними. В работе используется для
   моделирования последовательности фаз.
\item[\textbf{Сигнал агента}] Булев индикатор, выставляемый агентом
   в общем состоянии и наблюдаемый маршрутизатором. Примеры ---
   готовность к следующей фазе, признак тупикового состояния,
   счетчик отказов критика.
\item[\textbf{Маршрутизатор фаз} (\code{PhaseRouter})] Компонент,
   принимающий решение о переходе процесса в следующую фазу.
   Реализован в трех вариантах --- правиловом, на основе языковой
   модели и в виде оракула.
\item[\textbf{Маршрутизатор топологий} (\code{TopologyRouter})]
   Компонент, принимающий решение о смене активной базовой
   топологии внутри адаптивной мета-топологии.
\item[\textbf{Защитные правила переключения} (\code{SwitchGuards})]
   Набор ограничений, предотвращающих частое осцилляционное
   переключение топологий. Включает минимальное число итераций
   в одной топологии, интервал между сменами, лимит переключений
   за фазу и за весь запуск.
\item[\textbf{Ворота передачи состояния} (\code{TransitionGate})]
   Компонент, применяющий таблицу правил передачи общего
   состояния при смене активной топологии и регистрирующий факт
   перехода в журнале.
\end{description}

\subsection*{Человек в контуре управления}

\begin{description}[leftmargin=*, style=nextline]
\item[\textbf{Человек в контуре управления} (от англ. Human-in-the-
   Loop, HITL)] Парадигма, в которой автоматизированная система
   обращается к человеку за решением, оценкой или подтверждением
   в определенных точках процесса.
\item[\textbf{Шлюз взаимодействия с человеком} (\code{HumanGateway})]
   Программный протокол, через который агенты системы обращаются
   к человеку. Обеспечивает идемпотентность повторных вызовов по
   паре идентификаторов запуска и запроса.
\item[\textbf{Роль координатора}] Человек управляет планом и
   делегированием задач агентам. Активна в иерархических и
   полносвязных топологиях.
\item[\textbf{Роль рецензента}] Человек проверяет промежуточные
   результаты и инициирует их доработку. Естественно вписывается
   в Star и Chain.
\item[\textbf{Роль судьи}] Человек выносит финальный вердикт при
   наличии нескольких альтернативных решений. Применима в Debate.
\item[\textbf{Роль равноправного участника} (от англ. peer)]
   Человек участвует в обсуждении наравне с агентами, внося
   собственные сообщения в общую шину. Применима в Mesh.
\item[\textbf{Роль наблюдателя} (от англ. monitor)] Человек
   пассивно следит за процессом и вмешивается только в критических
   случаях --- превышение бюджета, зависание, циклы. Применима
   в любой топологии.
\item[\textbf{Маршрутизатор ролей} (\code{RoleRouter})] Компонент,
   динамически выбирающий активную роль человека в зависимости
   от фазы и наблюдаемого состояния.
\item[\textbf{Шкала NASA-TLX} (от англ. NASA Task Load Index)]
   Стандартный инструмент субъективной оценки когнитивной
   нагрузки по шести шкалам --- умственное, физическое и временное
   усилие, результативность, общее усилие, фрустрация.
\end{description}

\subsection*{Метрики}

\begin{description}[leftmargin=*, style=nextline]
\item[\textbf{Качество решения} ($q$)] Агрегированный показатель
   качества в диапазоне $[0, 1]$. Конкретная формула зависит от
   класса задач --- доля прошедших юнит-тестов, точное совпадение
   числа, среднее ROUGE-L с покрытием концептов.
\item[\textbf{Полная стоимость прогона} ($c$)] Сумма стоимостей
   всех вызовов языковых моделей за один прогон, выраженная в
   долларах США.
\item[\textbf{Стенное время прогона} ($\tau$)] Полное время
   выполнения одного прогона в секундах с момента старта до
   получения финального ответа.
\item[\textbf{Доля отказов} ($r_{\text{fail}}$)] Доля прогонов,
   завершившихся со статусом, отличным от \enquote{выполнено}.
   Включает превышение бюджета, превышение максимального числа
   итераций, ошибки вызова инструментов.
\item[\textbf{Число переключений топологии} ($N_{\text{switch}}$)]
   Фактическое число смен активной базовой топологии за один
   прогон адаптивной мета-топологии.
\item[\textbf{Доля заблокированных решений маршрутизатора}
   ($r_{\text{guard}}$)] Доля решений маршрутизатора топологий,
   заблокированных защитными правилами. Высокое значение
   свидетельствует о склонности маршрутизатора к частым
   переключениям.
\item[\textbf{Стоимостный вклад маршрутизатора} ($\gamma$)] Доля
   общего бюджета прогона, потраченная на собственные вызовы
   языковой модели маршрутизатором. Критическая метрика, отвечающая
   на вопрос, не превышает ли стоимость адаптации получаемый
   от нее выигрыш.
\item[\textbf{Доля регрессий качества} ($r_{\text{hurt}}$)] Доля
   задач, на которых адаптивная конфигурация дает качество хуже
   лучшей статической конфигурации. Ключевой индикатор
   безопасности механизма адаптации.
\item[\textbf{Разрыв до оракула} ($\Delta_{\text{LOO}}$)] Разность
   между качеством, обеспечиваемым оракулом по схеме исключения
   по одному, и качеством фактического маршрутизатора. Служит
   верхней границей достижимого качества при идеальной
   маршрутизации в рамках того же пространства решений.
\item[\textbf{Прокси-метрика когнитивной нагрузки}
   ($L_{\text{cog}}$)] Агрегированный показатель когнитивной
   нагрузки симулятора человека, объединяющий число обращений,
   объем предоставляемого контекста и среднюю задержку ответа.
   В экспериментах с реальными участниками заменяется на шкалу
   NASA-TLX.
\end{description}

\subsection*{Инфраструктура реализации}

\begin{description}[leftmargin=*, style=nextline]
\item[\textbf{LangGraph}] Библиотека на языке Python для построения
   и исполнения направленных графов состояния с поддержкой
   асинхронных вызовов, чекпойнтера и параллельного выполнения
   узлов.
\item[\textbf{Чекпойнтер}] Механизм сохранения промежуточного
   состояния графа в постоянное хранилище, позволяющий возобновлять
   прерванные запуски с последней успешной точки.
\item[\textbf{Apache Parquet}] Колоночный формат хранения больших
   объемов структурированных данных. В работе используется
   для журналов взаимодействия агентов и журналов вызовов
   языковых моделей.
\item[\textbf{Грид-эксперимент} (от англ. grid sweep)] Запуск
   серии прогонов по декартову произведению значений нескольких
   параметров с параллельным исполнением ячеек.
\item[\textbf{Прогон} (от англ. run)] Один полный цикл выполнения
   мульти-агентной системы от поступления задачи до получения
   ответа. Является элементарной единицей измерения в экспериментах.
\end{description}

\subsection*{Эксперимент и статистика}

\begin{description}[leftmargin=*, style=nextline]
\item[\textbf{Корпус задач}] Размеченное множество тестовых задач
   одного класса с заранее определенной процедурой автоматической
   оценки ответа. В работе используются четыре корпуса.
\item[\textbf{Воронка экспериментов}] Последовательность
   экспериментов, в которой каждый следующий сужает пространство
   рассматриваемых конфигураций по результатам предыдущего, что
   снижает суммарный объем вычислений без потери информативности.
\item[\textbf{Подтверждающий прогон} (от англ. confirmation run)]
   Дополнительный эксперимент на отличной от основной языковой
   модели, проводимый для проверки независимости выводов от
   выбора семейства весов.
\item[\textbf{Дисперсионный анализ} (от англ. analysis of
   variance, ANOVA)] Статистический метод проверки гипотезы о
   равенстве средних значений в нескольких группах путем
   сравнения межгрупповой и внутригрупповой дисперсии.
\item[\textbf{Поправка Бенджамини---Хохберга}] Процедура контроля
   доли ложных открытий при множественном статистическом
   сравнении.
\item[\textbf{Коэффициент эффекта Коэна} ($d$)] Стандартизированная
   мера величины различия между двумя средними. Используется
   для оценки практической значимости различий, а не только
   статистической.
\item[\textbf{Метод бутстрапа}] Непараметрический метод оценивания
   распределения статистики путем многократного ресэмплирования
   исходной выборки с возвращением. В работе применяется
   для построения доверительных интервалов средних значений
   метрик качества с уровнем доверия 95 процентов.
\item[\textbf{Доверительный интервал}] Интервал значений,
   накрывающий истинное значение параметра распределения с
   заданной вероятностью. В работе строится методом бутстрапа
   с уровнем доверия 95 процентов и тысячей повторений
   ресэмплирования.
\item[\textbf{Конфигурация системы}] Конкретный набор
   значений всех переменных эксперимента --- активная топология,
   роль человека, маршрутизаторы, языковая модель, корпус задач
   и значения сидов. Различают статические конфигурации (с
   фиксированной топологией) и адаптивные (с переключением
   топологии по фазам).
\item[\textbf{Метод исключения по одному} (от англ. leave-one-out,
   LOO)] Процедура построения верхнего потолка адаптации, в
   которой для каждой задачи оптимальная конфигурация вычисляется
   на остальных задачах того же класса.
\end{description}

\newpage
% =====================================================================
% Приложение В. Состав программного репозитория.
% Параграфный формат для компактности.
% =====================================================================

\section{Состав программного репозитория}
\label{app:repo}

В приложении приведено краткое описание состава программного
репозитория проекта. Состав фиксирован по состоянию идентификатора
фиксации, указанного в приложении~\ref{app:github}.

\begingroup
\small

\subsection*{Корневая структура}

\noindent\textbf{\code{src/atm/}} --- основной пакет фреймворка,
тринадцать функциональных слоев.\quad
\textbf{\code{conf/}} --- декларативная конфигурация экспериментов
и компонентов в формате YAML.\quad
\textbf{\code{scripts/}} --- самостоятельные скрипты обслуживания
экспериментов и сборки производных артефактов.\quad
\textbf{\code{notebooks/}} --- шаблон Jupyter-ноутбука для
аналитической постобработки результатов прогонов.\quad
\textbf{\code{pyproject.toml}} --- декларация пакета, зависимостей
и точки входа \code{atm}.\quad
\textbf{\code{README.md}} --- точка входа в документацию проекта
и инструкция по сборке среды.

\subsection*{Слои пакета \code{src/atm/}}

\noindent\textbf{\code{atm.core}} --- контейнеры состояния и
базовые типы (\code{AgentState}, \code{SharedState},
\code{GraphState}, \code{Phase}, \code{HumanRole}).\quad
\textbf{\code{atm.llm}} --- оболочка над провайдерами языковых
моделей, ценообразование, бюджетные ограничители,
детерминированный \code{FakeLLM}, фабрика \code{build\_llm}.\quad
\textbf{\code{atm.tools}} --- реестр инструментов и обертка
над Docker-песочницей.\quad
\textbf{\code{atm.agents}} --- базовый класс \code{Agent} и
пять специализированных ролей плюс координатор.\quad
\textbf{\code{atm.topology}} --- шесть реализаций топологий
(\code{star.py}, \code{chain.py}, \code{mesh.py}, \code{debate.py},
\code{hierarchical.py}, \code{adaptive.py}) поверх общего протокола
\code{Topology}.\quad
\textbf{\code{atm.phases}} --- конечный автомат фаз, маршрутизатор
топологий, защитные правила переключения, ворота передачи
состояния.\quad
\textbf{\code{atm.human}} --- протокол \code{HumanGateway},
реализации шлюзов, маршрутизатор ролей, политика тайм-аутов.\quad
\textbf{\code{atm.storage}} --- объектно-реляционное отображение,
модели таблиц PostgreSQL, асинхронные сессии, писатель Parquet,
обертка над чекпойнтером LangGraph, миграции \code{alembic/}.\quad
\textbf{\code{atm.observability}} --- перехватчик событий
LangGraph, сериализация, обертки над логгером.\quad
\textbf{\code{atm.tasks}} --- загрузчики и оценщики четырех
корпусов (HumanEval, GSM8K, CommonGen, InfiAgent-DABench).\quad
\textbf{\code{atm.evaluation}} --- судья оценки качества, оракулы
достоверности, агрегатор шкалы NASA-TLX.\quad
\textbf{\code{atm.experiment}} --- схемы конфигов на Pydantic,
загрузчик OmegaConf с расширением по сетке, командная строка
\code{atm} на Typer.\quad
\textbf{\code{atm.analysis}} --- загрузчики результатов,
агрегаторы метрик, генератор графиков, построитель оракула.

\subsection*{Конфигурационные файлы \code{conf/}}

\noindent\textbf{\code{conf/experiments/}} --- полные конфигурации
экспериментов E1--E4, подтверждающих прогонов и проверок
работоспособности.\quad
\textbf{\code{conf/agents/}} --- системные подсказки шести ролей
агентов.\quad
\textbf{\code{conf/topology/}} --- параметры по умолчанию для
шести топологий.\quad
\textbf{\code{conf/human/}} --- параметры шлюза HITL и системные
подсказки симулятора.\quad
\textbf{\code{conf/model/}} --- параметры языковых моделей.\quad
\textbf{\code{conf/task/}, \code{conf/tasks/}} --- загрузчики
корпусов и стратифицированные выборки.\quad
\textbf{\code{conf/oracle/}} --- конфигурации построения оракула.\quad
\textbf{\code{conf/evaluation/}} --- параметры судьи и агрегаторов.\quad
\textbf{\code{conf/sandbox/}} --- параметры Docker-песочницы и
профиль системных вызовов.

\subsection*{Скрипты и ноутбуки}

\noindent\textbf{\code{scripts/gen\_analysis\_notebook.py}} ---
генератор аналитического ноутбука из шаблона.\quad
\textbf{\code{scripts/derive\_seccomp.py}} --- сборка профиля
системных вызовов для песочницы.\quad
\textbf{\code{notebooks/analysis\_template.ipynb}} --- шаблонный
ноутбук из шести секций для постобработки и визуализации
результатов прогона.

\endgroup

\newpage
% =====================================================================
% Приложение Г. Открытый исходный код (четвертое приложение работы,
% буквенная нумерация А, Б, В, Г из main.tex).
% Ссылка на репозиторий + согласие на публикацию ВКР
% (требование методички, раздел 5.5 и 5.6.3).
% =====================================================================

\clearpage
\section{Открытый исходный код}
\label{app:github}

В соответствии с пунктом 5.5 Правил подготовки и защиты ВКР
полная программная реализация настоящего исследования размещена
в открытом доступе на платформе GitHub.

\begin{description}[leftmargin=2em, itemsep=0pt]
\item[\textbf{Адрес репозитория}]
   \url{https://github.com/cactus-tim/adaptive-topologies-mas}
\item[\textbf{Идентификатор фиксации}] коммит
   \code{058271a39a0e} (полный хеш ---
   \mbox{058271a39a0e4ecc0e3b47bd1d826c03fc33872b}) ветки
   \code{main}.
\item[\textbf{Лицензия}] MIT (файл \code{LICENSE}).
\item[\textbf{Сборка и воспроизведение}] инструкция в
   \code{README.md}; конфигурации экспериментов главы~\ref{sec:experiments}
   --- в \code{conf/experiments/}, запуск через
   \code{atm run~/~atm grid --config conf/experiments/e1.yaml}.
\end{description}

Согласие автора на публикацию ВКР на портале НИУ ВШЭ (п. 5.6.3
Правил) оформлено QR-кодом системы «LMS-Антиплагиат» в составе
пакета документов в ГЭК.


\end{document}
